% Options for packages loaded elsewhere
% Options for packages loaded elsewhere
\PassOptionsToPackage{unicode}{hyperref}
\PassOptionsToPackage{hyphens}{url}
\PassOptionsToPackage{dvipsnames,svgnames,x11names}{xcolor}
%
\documentclass[
  american,
  12pt]{article}
\usepackage{xcolor}
\usepackage{amsmath,amssymb}
\setcounter{secnumdepth}{5}
\usepackage{iftex}
\ifPDFTeX
  \usepackage[T1]{fontenc}
  \usepackage[utf8]{inputenc}
  \usepackage{textcomp} % provide euro and other symbols
\else % if luatex or xetex
  \usepackage{unicode-math} % this also loads fontspec
  \defaultfontfeatures{Scale=MatchLowercase}
  \defaultfontfeatures[\rmfamily]{Ligatures=TeX,Scale=1}
\fi
\usepackage{lmodern}
\ifPDFTeX\else
  % xetex/luatex font selection
\fi
% Use upquote if available, for straight quotes in verbatim environments
\IfFileExists{upquote.sty}{\usepackage{upquote}}{}
\IfFileExists{microtype.sty}{% use microtype if available
  \usepackage[]{microtype}
  \UseMicrotypeSet[protrusion]{basicmath} % disable protrusion for tt fonts
}{}
\makeatletter
\@ifundefined{KOMAClassName}{% if non-KOMA class
  \IfFileExists{parskip.sty}{%
    \usepackage{parskip}
  }{% else
    \setlength{\parindent}{0pt}
    \setlength{\parskip}{6pt plus 2pt minus 1pt}}
}{% if KOMA class
  \KOMAoptions{parskip=half}}
\makeatother
% Make \paragraph and \subparagraph free-standing
\makeatletter
\ifx\paragraph\undefined\else
  \let\oldparagraph\paragraph
  \renewcommand{\paragraph}{
    \@ifstar
      \xxxParagraphStar
      \xxxParagraphNoStar
  }
  \newcommand{\xxxParagraphStar}[1]{\oldparagraph*{#1}\mbox{}}
  \newcommand{\xxxParagraphNoStar}[1]{\oldparagraph{#1}\mbox{}}
\fi
\ifx\subparagraph\undefined\else
  \let\oldsubparagraph\subparagraph
  \renewcommand{\subparagraph}{
    \@ifstar
      \xxxSubParagraphStar
      \xxxSubParagraphNoStar
  }
  \newcommand{\xxxSubParagraphStar}[1]{\oldsubparagraph*{#1}\mbox{}}
  \newcommand{\xxxSubParagraphNoStar}[1]{\oldsubparagraph{#1}\mbox{}}
\fi
\makeatother


\usepackage{longtable,booktabs,array}
\usepackage{calc} % for calculating minipage widths
% Correct order of tables after \paragraph or \subparagraph
\usepackage{etoolbox}
\makeatletter
\patchcmd\longtable{\par}{\if@noskipsec\mbox{}\fi\par}{}{}
\makeatother
% Allow footnotes in longtable head/foot
\IfFileExists{footnotehyper.sty}{\usepackage{footnotehyper}}{\usepackage{footnote}}
\makesavenoteenv{longtable}
\usepackage{graphicx}
\makeatletter
\newsavebox\pandoc@box
\newcommand*\pandocbounded[1]{% scales image to fit in text height/width
  \sbox\pandoc@box{#1}%
  \Gscale@div\@tempa{\textheight}{\dimexpr\ht\pandoc@box+\dp\pandoc@box\relax}%
  \Gscale@div\@tempb{\linewidth}{\wd\pandoc@box}%
  \ifdim\@tempb\p@<\@tempa\p@\let\@tempa\@tempb\fi% select the smaller of both
  \ifdim\@tempa\p@<\p@\scalebox{\@tempa}{\usebox\pandoc@box}%
  \else\usebox{\pandoc@box}%
  \fi%
}
% Set default figure placement to htbp
\def\fps@figure{htbp}
\makeatother



\ifLuaTeX
\usepackage[bidi=basic]{babel}
\else
\usepackage[bidi=default]{babel}
\fi
% get rid of language-specific shorthands (see #6817):
\let\LanguageShortHands\languageshorthands
\def\languageshorthands#1{}
\ifLuaTeX
  \usepackage[english]{selnolig} % disable illegal ligatures
\fi


\setlength{\emergencystretch}{3em} % prevent overfull lines

\providecommand{\tightlist}{%
  \setlength{\itemsep}{0pt}\setlength{\parskip}{0pt}}



 
\usepackage[]{natbib}
\bibliographystyle{agsm}


\addtolength{\oddsidemargin}{-.5in}%
\addtolength{\evensidemargin}{-1in}%
\addtolength{\textwidth}{1in}%
\addtolength{\textheight}{1.7in}%
\addtolength{\topmargin}{-1in}%
\usepackage{tikz}
\usetikzlibrary{decorations.pathreplacing}
\usetikzlibrary{positioning}
\usetikzlibrary{shapes.arrows,shapes.symbols,arrows.meta}
\usepackage[no-weekday]{eukdate}
\usepackage{bm}
\usepackage{mathptmx}
\usepackage{amsmath} % Ensure amsmath is included for \eqref
\usepackage{bbm}
\usepackage[]{natbib}
%\usepackage[top=2cm, bottom=2cm, left=2.5cm, right=2.5cm]{geometry}
%\usepackage[letterpaper, margin=1in]{geometry}
\allowdisplaybreaks
\sloppy
\setlength{\parskip}{0pt}
\setlength{\parindent}{15pt}
\expandafter\def\expandafter\normalsize\expandafter{%
  \normalsize
  \setlength\abovedisplayskip{6pt}
  \setlength\belowdisplayskip{6pt}
  \setlength\abovedisplayshortskip{3pt}
  \setlength\belowdisplayshortskip{3pt}
}
\setlength{\bibsep}{0pt plus 0.3ex}
\usepackage{microtype}
\usepackage[onlyrefs]{mathtools}
\setlength{\tabcolsep}{0.11cm}
\usepackage{threeparttable}
\makeatletter
\g@addto@macro\TPT@defaults{\linespread{1}\selectfont}
\def\fps@table{!tbh}
\def\fps@figure{!tbh}
\makeatother

%% CAPTIONS
\usepackage{setspace}
\usepackage{caption}
\DeclareCaptionStyle{italic}[justification=centering]
 {labelfont={bf},textfont={it},labelsep=colon}
\captionsetup[figure]{style=italic,format=hang,singlelinecheck=true,font=singlespacing}
\captionsetup[table]{style=italic,format=hang,singlelinecheck=true,font=singlespacing}

\usepackage{titlesec}
\titlespacing*{\section}
  {0pt}{0.3\baselineskip}{0.1\baselineskip}
\titlespacing*{\subsection}
  {0pt}{0.3\baselineskip}{0.1\baselineskip}
\makeatletter
\@ifpackageloaded{caption}{}{\usepackage{caption}}
\AtBeginDocument{%
\ifdefined\contentsname
  \renewcommand*\contentsname{Table of contents}
\else
  \newcommand\contentsname{Table of contents}
\fi
\ifdefined\listfigurename
  \renewcommand*\listfigurename{List of Figures}
\else
  \newcommand\listfigurename{List of Figures}
\fi
\ifdefined\listtablename
  \renewcommand*\listtablename{List of Tables}
\else
  \newcommand\listtablename{List of Tables}
\fi
\ifdefined\figurename
  \renewcommand*\figurename{Figure}
\else
  \newcommand\figurename{Figure}
\fi
\ifdefined\tablename
  \renewcommand*\tablename{Table}
\else
  \newcommand\tablename{Table}
\fi
}
\@ifpackageloaded{float}{}{\usepackage{float}}
\floatstyle{ruled}
\@ifundefined{c@chapter}{\newfloat{codelisting}{h}{lop}}{\newfloat{codelisting}{h}{lop}[chapter]}
\floatname{codelisting}{Listing}
\newcommand*\listoflistings{\listof{codelisting}{List of Listings}}
\usepackage{amsthm}
\theoremstyle{plain}
\newtheorem{proposition}{Proposition}[section]
\theoremstyle{plain}
\newtheorem{corollary}{Corollary}[section]
\theoremstyle{remark}
\AtBeginDocument{\renewcommand*{\proofname}{Proof}}
\newtheorem*{remark}{Remark}
\newtheorem*{solution}{Solution}
\newtheorem{refremark}{Remark}[section]
\newtheorem{refsolution}{Solution}[section]
\makeatother
\makeatletter
\makeatother
\makeatletter
\@ifpackageloaded{caption}{}{\usepackage{caption}}
\@ifpackageloaded{subcaption}{}{\usepackage{subcaption}}
\makeatother
\usepackage{bookmark}
\IfFileExists{xurl.sty}{\usepackage{xurl}}{} % add URL line breaks if available
\urlstyle{same}
\hypersetup{
  pdftitle={Online conformal inference for multi-step time series forecasting},
  pdfauthor={Xiaoqian Wang; Rob J Hyndman},
  pdflang={en-US},
  pdfkeywords={Conformal prediction, Coverage
guarantee, Distribution-free inference, Exchangeability, Weighted
quantile estimate},
  colorlinks=true,
  linkcolor={blue},
  filecolor={Maroon},
  citecolor={Blue},
  urlcolor={Blue},
  pdfcreator={LaTeX via pandoc}}



\newcommand{\CPTSadd}[1]{\ifmmode\else\leavevmode\fi\pdfcolorstack 0 push{0 0 1 rg 0 0 1 RG}#1\pdfcolorstack 0 pop}
\begin{document}
\def\spacingset#1{\renewcommand{\baselinestretch}%
{#1}\small\normalsize} \spacingset{1}

\renewcommand*{\arraystretch}{0.5} % Specify row height in a table globally

%%%%%%%%%%%%%%%%%%%%%%%%%%%%%%%%%%%%%%%%%%%%%%%%%%%%%%%%%%%%%%%%%%%%%%%%%%%%%%

\title{\bf Online conformal inference for multi-step time series
forecasting}
\author{
Xiaoqian Wang\thanks{Corresponding author. Xiaoqian Wang, Academy of
Mathematics and Systems Science, Chinese Academy of Sciences, Beijing,
100190, China. E-mail address:
\href{mailto:xiaoqian.wang@amss.ac.cn}{xiaoqian.wang@amss.ac.cn} (X.
Wang).} \vspace{0.2em}\\
Academy of Mathematics and Systems Science, Chinese Academy of
Sciences \vspace{0.2em}\\
and \vspace{0.2em}\\Rob J Hyndman \vspace{0.2em}\\
Department of Econometrics \& Business Statistics, Monash
University \vspace{0.2em}\\
}
\maketitle

\bigskip
\bigskip
\begin{abstract}
We \CPTSadd{study} the \CPTSadd{construction} of distribution-free prediction intervals for
multi-step time series forecasting\CPTSadd{. In this setting}, forecast errors \CPTSadd{at
different horizons are serially dependent, and feedback is delayed,
since} the \CPTSadd{realized error of an} \(h\)-step-ahead forecast \CPTSadd{becomes
available only after \(h\) periods}. \CPTSadd{We} propose \CPTSadd{autocorrelated multi-step
conformal prediction} (AcMCP), \CPTSadd{an online procedure that accommodates
delayed feedback and explicitly} incorporates \CPTSadd{the dependence among}
multi-step forecast errors \CPTSadd{into conformal calibration}. \CPTSadd{To characterize
this dependence}, we \CPTSadd{derive first-order representations of the optimal
forecast errors under a general nonstationary autoregressive
data-generating process}, \CPTSadd{together with explicit bounds on} the \CPTSadd{remainder
terms. AcMCP guarantees long-run} coverage \CPTSadd{of at least} the \CPTSadd{nominal level
at every} horizon, \CPTSadd{without requiring exchangeability or any assumption on
how the data distribution evolves over time}. \CPTSadd{We also establish a
finite-sample bound on the} empirical \CPTSadd{coverage shortfall}, \CPTSadd{which depends
explicitly on the forecast horizon} and \CPTSadd{the evaluation sample size.
Notably}, \CPTSadd{these guarantees do not rely on the accuracy of the linear
approximations. Simulations and real-data applications show} that AcMCP
\CPTSadd{adapts to time-varying dynamics and maintains rolling} coverage \CPTSadd{close to}
the \CPTSadd{nominal level. By exploiting the dependence among multi-step errors},
\CPTSadd{it often yields narrower} intervals \CPTSadd{than comparably well-calibrated
methods, particularly at longer horizons}.
\end{abstract}

\noindent%
{\it Keywords:} Conformal prediction, Coverage
guarantee, Distribution-free inference, Exchangeability, Weighted
quantile estimate
\vfill

\newpage
\spacingset{1.8} % DON'T change the spacing!

\renewcommand{\theproposition}{\arabic{proposition}} % numbering propositions sequentially and individually by type and not by section
\setcounter{proposition}{0}

\renewcommand{\thecorollary}{\arabic{corollary}} % numbering corollaries sequentially and individually
\counterwithout{corollary}{section}
\setcounter{corollary}{0}

\section{Introduction}\label{sec-intro}

Conformal prediction \citep{papadopoulos2002, vovk2005} is a
\CPTSadd{model-agnostic} framework for uncertainty quantification \CPTSadd{that} constructs
prediction intervals \CPTSadd{with valid} coverage\CPTSadd{,} without \CPTSadd{requiring} stringent
assumptions on the \CPTSadd{underlying} data distribution\CPTSadd{. Its classical
guarantees}, \CPTSadd{however, rely on} the data \CPTSadd{being} i.i.d. or, more generally,
exchangeable. \CPTSadd{In time series forecasting, observations are typically
serially dependent} and \CPTSadd{may exhibit nonstationarity} or \CPTSadd{distributional
shifts}, \CPTSadd{so exchangeability generally fails} and \CPTSadd{the standard guarantees
no longer hold}.

\CPTSadd{Recent work has sought} to \CPTSadd{relax this requirement}. \CPTSadd{One line} of \CPTSadd{research
replaces} exchangeability \CPTSadd{with weaker}, \CPTSadd{time-series-specific} assumptions\CPTSadd{.
To handle temporal dependence}, \CPTSadd{approximate} coverage guarantees \CPTSadd{have been
obtained through permutation schemes that preserve} the \CPTSadd{temporal
dependence structure \citep{chernozhukov2018, yu2022}} and through
\CPTSadd{concentration and decoupling arguments for split conformal} prediction
\CPTSadd{\citep{oliveira2024}}. \CPTSadd{\citet{xu2021} use bootstrap ensemble predictors
to obtain out-of-bag residuals and update their empirical distribution}
as \CPTSadd{new observations arrive}, \CPTSadd{establishing approximate coverage under
stationary}, \CPTSadd{strongly mixing errors}. \CPTSadd{Under} covariate shift,
\CPTSadd{\citet{tibshirani2019}} and \CPTSadd{\citet{yang2024} establish coverage
guarantees when} the \CPTSadd{training and test distributions differ only in their
covariate marginals}. \CPTSadd{Other work generalizes conformal calibration by
replacing the unweighted empirical distribution of} nonconformity scores
\CPTSadd{with a weighted one \citep{barber2023, guan2023}}.

\CPTSadd{A separate strand} of research \CPTSadd{develops online conformal} methods that
\CPTSadd{adjust prediction intervals using sequential} coverage \CPTSadd{feedback}. \CPTSadd{Adaptive}
conformal \CPTSadd{inference \citep[ACI,][]{gibbs2021} updates the working
miscoverage level in response to observed coverage errors}, \CPTSadd{so} that
\CPTSadd{empirical coverage converges} to the \CPTSadd{target level without distributional}
assumptions on the time series\CPTSadd{. Subsequent extensions use expert
aggregation} and \CPTSadd{adaptive online learning to reduce sensitivity to} the
\CPTSadd{learning rate} and to \CPTSadd{improve adaptation to} distribution \CPTSadd{shifts
\citep{zaffran2022, gibbs2024, bhatnagar2023}}. \CPTSadd{Conformal PID control
\citep{angelopoulos2024} builds} on the \CPTSadd{same feedback principle but
adjusts the thresholds} of \CPTSadd{the nonconformity scores directly}, \CPTSadd{combining
quantile tracking}, \CPTSadd{correction for cumulative} coverage \CPTSadd{errors}, \CPTSadd{and a
scorecaster that exploits predictable temporal structure in} the \CPTSadd{scores}.
\CPTSadd{Further refinements improve stability through decaying} step sizes
\CPTSadd{\citep{angelopoulos2024online}}, \CPTSadd{reduce tuning through parameter-free
updates \citep{podkopaev2024}}, and \CPTSadd{incorporate smooth feedback} that
\CPTSadd{reflects} the \CPTSadd{magnitude} of the \CPTSadd{discrepancy between scores} and \CPTSadd{thresholds
\citep{wu2025}}.

\CPTSadd{Most of this literature, however, targets} single-step forecasting,
\CPTSadd{whereas} many applications require uncertainty quantification over
\CPTSadd{several} future horizons. \CPTSadd{Multi-step forecasting poses two additional
challenges for} conformal \CPTSadd{methods. First}, \CPTSadd{feedback is delayed:} the
\CPTSadd{realized error of an \(h\)-step-ahead forecast becomes available only
\(h\) periods after the forecast is issued}, \CPTSadd{so online updates must rely
on outdated information. Second}, \CPTSadd{forecast errors} at different horizons
\CPTSadd{issued from the same origin} are \CPTSadd{dependent}, \CPTSadd{because errors propagate
through the forecasting recursion.}

\CPTSadd{Existing conformal methods for multi-step forecasting address these
challenges only partially.} \citet{stankeviciute2021} integrate conformal
prediction with recurrent neural networks for multi-step forecasting and
apply \CPTSadd{a} Bonferroni correction to \CPTSadd{obtain simultaneous} coverage \CPTSadd{across
horizons}, \CPTSadd{but their guarantee relies on exchangeability across} time
series. \CPTSadd{\citet{schlembach2025} extend this approach to multivariate
targets under nonexchangeability, with a joint miscoverage bound across
horizons and response dimensions.} Bellman conformal inference
\CPTSadd{\citep{yang2024ts} uses} multi-step \CPTSadd{forecasts in} a \CPTSadd{look-ahead
optimization} that \CPTSadd{trades off} interval length \CPTSadd{against a} miscoverage
\CPTSadd{penalty}, \CPTSadd{but it outputs, and guarantees coverage for, only the resulting
sequence of one-step-ahead intervals; it also requires nested interval
families,} which \CPTSadd{calls for additional construction when only point
forecasts are available,} and \CPTSadd{solves a dynamic program} at \CPTSadd{every} time
step. \CPTSadd{The multi-step extension of sequential predictive conformal
inference \citep[SPCI,][]{xu2023} fits a separate predictor and
conditional quantile model for each horizon and uses past residuals} to
\CPTSadd{adapt the interval bounds}; \CPTSadd{although SPCI has asymptotic conditional
coverage guarantees under suitable assumptions}, \CPTSadd{its multi-step
implementation is not analyzed separately. To our knowledge, existing
multi-step conformal methods seldom provide coverage guarantees without
exchangeability, and none exploits the dependence among forecast errors
across horizons. A method that handles delayed feedback, accounts} for
\CPTSadd{this dependence}, \CPTSadd{and retains distribution-free coverage guarantees is
therefore still lacking}.

We consider a time series \CPTSadd{\(\{y_t\}_{t\ge1}\)} generated by an unknown\CPTSadd{,
possibly nonstationary data-generating} process (DGP) \CPTSadd{that} may depend on
its own past \CPTSadd{and on} exogenous predictors \(\bm{x}_t\in\mathbb{R}^p\). At
each time \(t\), we \CPTSadd{wish} to \CPTSadd{construct} a prediction interval
\(\hat{\mathcal{C}}_{t+h|t}\) for \(y_{t+h}\) \CPTSadd{at every horizon}
\(h\in[H]\). Given a \CPTSadd{user-specified miscoverage level
\(\alpha\in(0,1)\)}, the goal is \CPTSadd{horizon-specific} long-run coverage\CPTSadd{,
namely
\(\liminf_{T\to\infty}\frac{1}{T}\sum_{t=1}^{T}\mathbbm{1}\left\{y_{t+h}\in\hat{\mathcal{C}}_{t+h|t}\right\}\geq 1-\alpha\),}
for \CPTSadd{\(h\in[H]\)}.

\CPTSadd{This paper makes three contributions. First, we develop multi-step}
conformal prediction \CPTSadd{methods within} an online learning \CPTSadd{framework with}
sequential \CPTSadd{training and calibration splits}. \CPTSadd{Building on conformal PID
control of \citet{angelopoulos2024}}, we \CPTSadd{introduce its} multi-step
\CPTSadd{extension (MPID)}, \CPTSadd{which accommodates delayed} coverage \CPTSadd{feedback. Second},
\CPTSadd{we propose autocorrelated} multi-step conformal \CPTSadd{prediction} (\CPTSadd{AcMCP}), which
\CPTSadd{explicitly incorporates the dependence among multi-step forecast errors
into the scorecaster}. \CPTSadd{To motivate its design}, we \CPTSadd{analyze optimal
forecast errors} under a general nonstationary autoregressive DGP \CPTSadd{with
future exogenous predictors treated as known}, \CPTSadd{derive first-order error
decompositions, and bound} the \CPTSadd{Taylor remainders under suitable
smoothness and moment conditions. These results lead to} a \CPTSadd{scorecaster
that combines a moving average model for errors} at \CPTSadd{a fixed horizon} with
\CPTSadd{a regression on errors across horizons from the same} forecast \CPTSadd{origin.
AcMCP combines these error forecasts with online calibration of the two
tails to construct asymmetric prediction intervals for each} horizon.
Third, \CPTSadd{for both MPID and} AcMCP, \CPTSadd{we establish finite-sample bounds on
empirical coverage that} explicitly account for \CPTSadd{delayed feedback and
imply long-run coverage} of at \CPTSadd{least} the \CPTSadd{nominal level at each fixed
horizon. These guarantees hold under suitable conditions on} the scores
\CPTSadd{and} the \CPTSadd{calibration procedure}, without any assumption on \CPTSadd{how the} data
distribution \CPTSadd{changes over time}. \CPTSadd{They concern coverage averaged over time
at} each horizon, \CPTSadd{rather than simultaneous coverage across horizons or
conditional coverage at individual forecast origins}. \CPTSadd{Importantly}, the
\CPTSadd{coverage guarantees} of \CPTSadd{AcMCP do not depend on the accuracy of the
forecast-error approximations; the contribution of the scorecaster to
interval efficiency is assessed empirically. Simulations} and
applications to electricity demand and eating-out expenditure show that
AcMCP adapts to \CPTSadd{distributional} changes and \CPTSadd{often produces} narrower
intervals \CPTSadd{than competing methods with comparable coverage}, particularly
\CPTSadd{at longer} horizons.

We developed the \texttt{conformalForecast} package for R to implement
the proposed multi-step conformal prediction methods\CPTSadd{. The} package is
publicly available on CRAN. All the data and code to reproduce the
experiments are made available at \url{https://github.com/xqnwang/cpts}.

\section{Online learning with sequential splits}\label{sec-setup}

\textbf{\CPTSadd{Problem setup.}} \CPTSadd{We observe a time series \(\{y_t\}_{t\geq1}\)
generated by an unknown DGP, which may depend on its own past as well as
on exogenous predictors
\(\bm{x}_t=(x_{1,t},\dots,x_{p,t})^{\prime}\in\mathbb{R}^p\) and their
histories. Write \(z_t=(\bm{x}_t,y_t)\). The joint distribution of
\(\{z_t\}_{t\geq1}\) is allowed to vary over time, thereby accommodating
nonstationary processes, and we impose no assumption on how it evolves.
At each time \(t\), we forecast \(H\) steps ahead and, for each
\(h\in[H]\), construct a} \emph{\CPTSadd{prediction set}}
\CPTSadd{\(\hat{\mathcal{C}}^{\alpha}_{t+h|t}\) for \(y_{t+h}\), which is a
prediction interval in this setting. The \(h\)-step-ahead forecast uses}
the \CPTSadd{observed} data \CPTSadd{\(\{z_i\}_{1\leq i\leq t}\) together with} the
\CPTSadd{exogenous} predictors \CPTSadd{\(\{\bm{x}_{t+j}\}_{1\leq j\leq h}\) over the
forecast period. These future predictors may be replaced by forecasts
based on information available up to and including} time \(t\)
\CPTSadd{(}\emph{\CPTSadd{ex-ante}} \CPTSadd{forecasts), or the actual values may be assumed
available (}\emph{\CPTSadd{ex-post}} \CPTSadd{forecasts)}. \CPTSadd{Given a user-specified}
\emph{\CPTSadd{miscoverage level}} \CPTSadd{\(\alpha\in(0,1)\), we seek intervals with
horizon-specific long-run coverage,}
\begin{equation}\CPTSadd{\phantomsection\label{eq-longrun}{
\liminf_{T\to\infty}\frac{1}{T-h}\sum_{t=h+1}^{T}\mathbbm{1}\left\{y_t\in\hat{\mathcal{C}}^{\alpha}_{t|t-h}\right\}\geq 1-\alpha,\qquad h\in[H].
}}\end{equation} \CPTSadd{This criterion concerns coverage averaged over time at
each forecast horizon, rather than coverage conditional on the available
information or simultaneous coverage of the entire future path. The
parametric form of the DGP is assumed only for the forecast-error
analysis in Section~\ref{sec-ppt}; the coverage guarantees in
Section~\ref{sec-cov} do not require it.}

\textbf{\CPTSadd{Forecasts and errors.}} Suppose that, at each time \(t\), a
forecasting model \(\hat{f}_t\) is trained on historical data
\(z_{1:t}\) and updated (retrained) as new observations arrive. \CPTSadd{Using
\(\hat{f}_t\), we produce point forecasts \(\hat{y}_{t+h|t}\) and signed
forecast errors \(e_{t+h|t}=y_{t+h}-\hat{y}_{t+h|t}\) for \(h\in[H]\).
Conformal inference then calibrates these errors to turn} the \CPTSadd{point
forecasts} into \CPTSadd{prediction intervals.}

\textbf{\CPTSadd{Prediction intervals.}} \CPTSadd{We construct generally asymmetric
intervals by calibrating} the \CPTSadd{lower and upper tails separately. Choose
tail levels \(\alpha^{-},\alpha^{+}>0\) satisfying
\(\alpha^{-}+\alpha^{+}=\alpha\). The two calibration procedures use the}
\emph{\CPTSadd{nonconformity scores}} \CPTSadd{\(s_{t+h|t}^{-}=-e_{t+h|t}\) and
\(s_{t+h|t}^{+}=e_{t+h|t}\)}, \CPTSadd{respectively}, \CPTSadd{and produce thresholds
\(q_{t+h|t}^{-}\) and \(q_{t+h|t}^{+}\)}. \CPTSadd{These thresholds define the
prediction interval}
\begin{equation}\CPTSadd{\phantomsection\label{eq-two-sided-pi}{
\hat{\mathcal{C}}_{t+h|t}^{\alpha}
:=\left[\hat y_{t+h|t}-q_{t+h|t}^{-},\ \hat y_{t+h|t}+q_{t+h|t}^{+}\right].
}}\end{equation} \CPTSadd{For each tail \(a\in\{-,+\}\)}, the \emph{\CPTSadd{miscoverage
indicator}}
\CPTSadd{\(\mathrm{err}_{t+h|t}^{a}:=\mathbbm{1}\{s_{t+h|t}^{a}>q_{t+h|t}^{a}\}\)
records whether} the \CPTSadd{observation falls beyond} the \CPTSadd{corresponding interval
bound}. Using \CPTSadd{opposite signed scores allows both tails to be calibrated
through} the \CPTSadd{same score-exceedance rule}, \CPTSadd{with observations at either
endpoint counted as covered. Because \(y_{t+h}\)}, \CPTSadd{and hence
\(s_{t+h|t}^{a}\) and \(\mathrm{err}_{t+h|t}^{a}\)}, \CPTSadd{is revealed only at
time \(t+h\), feedback is delayed: at time \(t\),} the \CPTSadd{information
available} for \CPTSadd{horizon \(h\) comes from forecast origins \(t-h\) and
earlier}. We \CPTSadd{use \(\alpha^{-}=\alpha^{+}=\alpha/2\) throughout} the
\CPTSadd{empirical analysis.}

\CPTSadd{Unless otherwise stated, we present} the \CPTSadd{methods and theoretical results
in subsequent sections for a generic single tail}, \CPTSadd{omitting the tail
superscript for simplicity. Accordingly}, \CPTSadd{\(\alpha\) denotes} the target
\CPTSadd{tail miscoverage} level \CPTSadd{in single-tail statements and the total
miscoverage level in two-sided statements}. For \CPTSadd{a candidate response
\(y\)}, \CPTSadd{\(s_{t+h|t}^{y}\) denotes} the \CPTSadd{score obtained by replacing
\(y_{t+h}\) with \(y\)}.

\textbf{Sequential split.} In a time series context, it is inappropriate
to perform \emph{random splitting}, a standard strategy in much of the
conformal prediction literature, due to the temporal dependency present
in the data. Following related work such as \citet{wisniewski2020} and
\citet{zaffran2022}, \CPTSadd{all} conformal methods proposed in this paper use a
\emph{sequential split} to preserve the temporal structure. For example,
the \(t\) available data points, \(z_{1:t}\), are sequentially split
into two consecutive sets, a \emph{proper training set}
\(\mathcal{D}_{\text{tr}} \subset \{1, \dots, t_r\}\) and a
\emph{calibration set}
\(\mathcal{D}_{\text{cal}} \subset \{t_r+1, \dots, t\}\), where
\(t_c=t-t_r \gg H\),
\(\mathcal{D}_{\text{tr}} \cup \mathcal{D}_{\text{cal}} = \{1, \dots, t\}\),
and
\(\mathcal{D}_{\text{tr}} \cap \mathcal{D}_{\text{cal}} = \varnothing\).
Here, ``proper'' means that the training set is used exclusively for
fitting the model, with no overlap \CPTSadd{with} the calibration set, which is
essential for the validity of coverage in conformal prediction
\citep{papadopoulos2002, vovk2005}. With sequential splitting, multiple
\(H\)-step forecasts and their nonconformity scores can be computed on
the calibration set.

\textbf{Online learning.} We \CPTSadd{adopt} the following generic online learning
framework, closely related to \citet{zaffran2022}, for all conformal
methods discussed in later sections. The procedure is illustrated in
\autoref{fig-flowchart}. This framework updates prediction intervals as
new data points arrive, allowing us to assess their long-run coverage
behavior. It adopts a standard rolling window evaluation strategy,
although expanding windows could easily be used instead.

\begin{figure}[!hbtp]
\vspace*{-1.5cm}
\hspace*{-2.5cm}
\centering
\scalebox{0.85}{
  \input figs/annotation.tex
}
\vspace*{-0.5cm}

\caption{Online learning framework with sequential splits. White: unused data; Gray: training data; Black: forecasts in \CPTSadd{the} calibration set; Blue: forecasts in \CPTSadd{the} test set.}\label{fig-flowchart}
\end{figure}

\begin{enumerate}
\def\labelenumi{\arabic{enumi}.}
\item
  \emph{Initialization}. Train a forecasting model on the initial proper
  training set \(z_{(1+t-t_r):t}\), setting \(t=t_r\). Then generate
  \(H\)-step-ahead point forecasts \(\{\hat{y}_{t+h|t}\}_{h\in[H]}\) and
  compute the corresponding \CPTSadd{upper- and lower-tail} nonconformity scores
  \CPTSadd{\(\{s_{t+h|t}^{+},s_{t+h|t}^{-}\}_{h\in[H]}\)} based on the true values
  \(\{y_{t+h}\}_{h\in[H]}\).
\item
  \emph{Recurring procedure}. Roll the training set forward by one data
  point by setting \(t \rightarrow t+1\). Then repeat step 1 until the
  nonconformity scores on the entire initial calibration set,
  \CPTSadd{\(\{s_{t+h|t}^{+},s_{t+h|t}^{-}\}_{t_r \leq t \leq t_r+t_c-h}\)} for
  \(h\in[H]\), are computed.
\item
  \emph{Quantile estimation and prediction interval calculation}. \CPTSadd{For
  each tail, use the} nonconformity scores from the calibration set to
  \CPTSadd{estimate horizon-specific quantiles. Combine the two one-sided
  thresholds via Equation \eqref{eq-two-sided-pi} to} compute
  \(H\)-step-ahead prediction intervals on the test set.
\item
  \emph{Online updating}. Continuously roll the training and calibration
  set\CPTSadd{s} forward, one data point at a time, to update the nonconformity
  scores \CPTSadd{used} for calibration, and repeat step 3 until prediction
  intervals for the entire test set are obtained. \CPTSadd{These intervals are}
  \(\{\hat{\mathcal{C}}_{t+h|t}^{\alpha}\}_{t_r+t_c \leq t \leq T-H}\)
  for \(h \in [H]\), where \(T-t_r-t_c\) is the length of the test set
  used to \CPTSadd{evaluate} coverage \CPTSadd{against the long-run target} in \CPTSadd{Equation
  \eqref{eq-longrun}}.
\end{enumerate}

\section{\CPTSadd{Multi-step} conformal \CPTSadd{PID control}}\label{sec-mpid}

\CPTSadd{For} a \CPTSadd{well-specified forecasting} model, the \CPTSadd{conditional} mean \CPTSadd{minimizes}
the expected squared forecast error, \CPTSadd{and} for \CPTSadd{covariance-stationary}
linear \CPTSadd{processes} the variance of the \CPTSadd{resulting} error \(e_{t+h|t}\) is
non-decreasing in \(h\) \citep{diebold1996, patton2007}. \CPTSadd{The
distribution} of nonconformity scores \CPTSadd{therefore differs across horizons},
\CPTSadd{so} the \CPTSadd{thresholds should be calibrated separately} for each \CPTSadd{\(h\in[H]\)}.
\CPTSadd{Calibration also need} to \CPTSadd{cope} with \CPTSadd{delayed feedback}: at time \(t\), the
\CPTSadd{realized} \(h\)-step \CPTSadd{score} is \CPTSadd{available only} for \CPTSadd{forecasts issued at
origin \(t-h\)} or \CPTSadd{earlier}.

We introduce \CPTSadd{the} \textbf{multi-step conformal PID control} method
(hereafter MPID), which extends the PID method \citep{angelopoulos2024},
originally developed for one-step-ahead forecasting\CPTSadd{, to the multi-step
setting within the online learning framework of Section~\ref{sec-setup}}.
For each forecast horizon \(h\in[H]\), \CPTSadd{define} the \CPTSadd{P-control state}
\begin{equation}\CPTSadd{\phantomsection\label{eq-p-state}{
q^{\mathrm P}_{t+h|t}=q^{\mathrm P}_{t+h-1|t-1}+\eta(\mathrm{err}_{t|t-h}-\alpha),
}}\end{equation} \CPTSadd{with a finite initial state \(q^{\mathrm P}_{2h|h}\),
typically zero. The} estimated \(1-\alpha\) quantile of the
\(h\)-step-ahead score at time \(t\) is updated as
\begin{equation}\CPTSadd{\phantomsection\label{eq-mpid}{
q_{t+h|t}=\underbrace{q^{\mathrm P}_{t+h|t}}_{\mathrm P}+\underbrace{r_t\Bigg(\sum_{i=h+1}^t(\mathrm{err}_{i|i-h}-\alpha)\Bigg)}_{\mathrm I}+\underbrace{\hat{s}_{t+h|t}}_{\mathrm D},
}}\end{equation} where \(\eta > 0\) is a constant learning rate, and
\(r_t\) is a saturation function that adheres to
\begin{equation}\phantomsection\label{eq-saturation_h}{
x \geq c \cdot g(t-h) \Longrightarrow r_t(x) \geq b, \quad \text {and} \quad x \leq-c \cdot g(t-h) \Longrightarrow r_t(x) \leq -b,
}\end{equation} for constant\CPTSadd{s} \(b, c > 0\) and an admissible function
\(g\) that is nonnegative, nondecreasing \CPTSadd{and sublinear}, \CPTSadd{i.e.,
\(g(n)/n\to0\)}. \CPTSadd{For} \(h=1\), Equation \eqref{eq-mpid} \CPTSadd{reduces} to the
\CPTSadd{standard} PID update. \CPTSadd{For each \(h\in[H]\)}, we \CPTSadd{apply this update
separately to the lower and upper tails}, \CPTSadd{using the tail levels
\(\alpha^{-}\) and \(\alpha^{+}\) and the feedback available at time
\(t\), and combine the resulting thresholds via Equation
\eqref{eq-two-sided-pi}} to \CPTSadd{form the prediction interval for \(y_{t+h}\)}.

The \CPTSadd{three components play distinct roles}.

\emph{\CPTSadd{P (quantile tracking).}} The \CPTSadd{P state moves up or down according} to
\CPTSadd{whether \(y_t\) fell outside} or \CPTSadd{inside} the \CPTSadd{one-sided} prediction set
\CPTSadd{issued} at time \CPTSadd{\(t-h\)}. \CPTSadd{It thus incorporates} the \CPTSadd{\(h\)-period feedback
delay explicitly:} the \CPTSadd{most recent feedback usable at time \(t\) is
\(\mathrm{err}_{t|t-h}\)}.

\emph{\CPTSadd{I (error integration).}} The I \CPTSadd{term accumulates the} historical
coverage errors \CPTSadd{of the} \(h\)-step-ahead \CPTSadd{one-sided} prediction \CPTSadd{sets,}
during updates, thereby enhancing the stability of coverage.

\emph{\CPTSadd{D (scorecasting).}} The D \CPTSadd{term} \(\hat{s}_{t+h|t}\) \CPTSadd{is a} forecast of
the \CPTSadd{\(h\)-step-ahead} nonconformity score, produced by any suitable
forecasting model (\CPTSadd{a} ``scorecaster''\CPTSadd{;} see Section~\ref{sec-acmcp})
trained \CPTSadd{on} the \(h\)-step-ahead scores available up to and including
time \(t\). \CPTSadd{Its} effectiveness depends on the ability of the scorecaster
to capture systematic and predictable \CPTSadd{structure} in the scores, such as
temporal dependence or conditional heteroskedasticity. Only when such
structure is present and adequately modeled \CPTSadd{does the D term provide}
additional benefits; otherwise, it \CPTSadd{can} increase the \CPTSadd{variability of}
coverage and \CPTSadd{of the} prediction \CPTSadd{sets}.

MPID \CPTSadd{is a special case of the generic update} in \CPTSadd{Equation
\eqref{eq-acmcp_1}, with
\(\hat q_{t+h|t}=q^{\mathrm P}_{t+h|t}+\hat s_{t+h|t}\). Its one- and
two-sided prediction sets therefore inherit} the \CPTSadd{finite-sample empirical
coverage bounds and long-run guarantees of Corollary~\ref{cor-cov_acmcp}
and Corollary~\ref{cor-cov-two-sided}}, \CPTSadd{respectively, under} the \CPTSadd{stated
assumptions. For reference, Appendix \ref{sec-extensions} presents
multi-step extensions} of \CPTSadd{split conformal prediction \citep{vovk2005},
conformal prediction} with \CPTSadd{weighted quantiles \citep{barber2023}}, \CPTSadd{and ACI
\citep{gibbs2021}; we do not establish} coverage guarantees for \CPTSadd{these}.

\section{Autocorrelated multi-step conformal
prediction}\label{sec-acmcp}

\CPTSadd{The scorecaster in MPID} is a model trained on the score sequence to
forecast the future score. The rationale behind it is to identify any
leftover signal in the score distribution not captured by the base
forecasting model. While this is appropriate in the context of huge data
sets and potentially weak learners, it is unlikely to be a useful
strategy for time series models. We expect to use a forecasting model
that leaves only white noise in the residuals (equivalent to the
one-step nonconformity scores). Moreover, the inclusion of a scorecaster
often only introduces variance to the quantile estimate, resulting in
inefficient (wider) prediction intervals.

On the other hand, in any non-trivial context, multi-step forecasts are
correlated with each other. Specifically, the \(h\)-step-ahead forecast
errors \(e_{t+h|t}\) are correlated with the forecast errors from the
previous \(h-1\) steps, as errors propagate through the recursive
forecasting procedure over the forecast horizon. However, no conformal
prediction methods have taken this potential dependence into account in
their methodological construction.

\CPTSadd{This} section \CPTSadd{studies} the \CPTSadd{structure} of multi-step forecast errors\CPTSadd{, uses
it to design the AcMCP scorecaster,} and \CPTSadd{establishes} coverage guarantees
\CPTSadd{that do not depend on how accurately the scorecaster approximates the
errors}.

\subsection{Properties of multi-step forecast errors}\label{sec-ppt}

We assume that a time series \(\{y_t\}_{t \geq 1}\) follows a general
nonstationary autoregressive process:
\begin{equation}\phantomsection\label{eq-dgp}{
y_t = f_{t}(y_{(t-d):(t-1)}, \bm{x}_{(t-k):t}) + \varepsilon_t,
}\end{equation} where \(f_{t}\) \CPTSadd{depends on} \(d\) lagged values of
\(y_t\) \CPTSadd{and on} the current \CPTSadd{and \(k\) preceding values} of the exogenous
predictors, and \(\varepsilon_t\) is \CPTSadd{the} innovation. \CPTSadd{To isolate
uncertainty in future responses, we treat future exogenous predictors as
known}. Let \CPTSadd{\(\mathcal F_t\)} be \CPTSadd{the information} generated by \CPTSadd{the
responses observed through time \(t\)} and the \CPTSadd{exogenous predictor path}.
The optimal \CPTSadd{point forecast under squared loss} is
\CPTSadd{\(\hat y_{t+h|t}=\mathbb E(y_{t+h}|\mathcal F_t)\)}, \CPTSadd{with error
\(e_{t+h|t}=y_{t+h}-\hat y_{t+h|t}\)}.

\CPTSadd{At origin \(t\)}, \CPTSadd{\(p_h=\min\{d,h-1\}\) of} the \CPTSadd{lagged responses entering
\(f_{t+h}\)} are \CPTSadd{unobserved}. \CPTSadd{For \(h\geq2\)}, \CPTSadd{collect them in
\(\bm z_{t,h}=(y_{t+h-1},\dots,y_{t+h-p_h})^{\prime}\)}, with \CPTSadd{conditional
mean \(\bm m_{t,h}=\mathbb E(\bm z_{t,h}|\mathcal F_t)\)}, \CPTSadd{and let
\(F_{t,h}\) denote \(f_{t+h}\) as a function of \(\bm z_{t,h}\) with all
other arguments held fixed}. \CPTSadd{We write \(\partial_i\) for differentiation
in the \(i\)th coordinate}, \CPTSadd{and \(\|\cdot\|_2\) and
\(\|\cdot\|_{\mathrm{op}}\) for} the \CPTSadd{Euclidean and operator norms},
\CPTSadd{respectively.}

\begin{proposition}[Autoregressive representation of optimal multi-step
forecast errors]\protect\hypertarget{prp-ar}{}\label{prp-ar}

\CPTSadd{Consider} the \CPTSadd{process in Equation \eqref{eq-dgp}}. \CPTSadd{Assume that each
\(f_s\) is twice continuously differentiable in} the \CPTSadd{lagged responses},
\CPTSadd{that \(\mathbb E(y_s^2)<\infty\) and
\(\mathbb E(\varepsilon_s|\mathcal F_{s-1})=0\) for all \(s\geq1\)}, \CPTSadd{and}
that the \CPTSadd{first-order Taylor terms and remainder below are integrable}.
\CPTSadd{Then}, \CPTSadd{for each \(h\in[H]\)},
\begin{equation}\CPTSadd{\phantomsection\label{eq-ar}{
e_{t+h|t}
=\varepsilon_{t+h}+\sum_{i=1}^{p_h}\phi_{t,h,i}e_{t+h-i|t}+\widetilde R_{t,h},
}}\end{equation} \CPTSadd{where \(\phi_{t,h,i}=\partial_iF_{t,h}(\bm m_{t,h})\)
and \(\widetilde R_{t,h}=R_{t,h}-\mathbb E(R_{t,h}|\mathcal F_t)\)}, \CPTSadd{with
\(R_{t,h}\)} the \CPTSadd{remainder} of the \CPTSadd{first-order Taylor expansion of
\(F_{t,h}\) at \(\bm m_{t,h}\)}. \CPTSadd{In particular},
\CPTSadd{\(\mathbb E(\widetilde R_{t,h}|\mathcal F_t)=0\)} and
\CPTSadd{\(\widetilde R_{t,1}=0\)}.

\CPTSadd{Moreover, for \(h\geq2\), if there is a finite deterministic
\(M_{t,h}\geq0\) with} \[
\CPTSadd{\sup_{u\in[0,1]}
\left\|\nabla^2F_{t,h}\!\left(\bm m_{t,h}+u(\bm z_{t,h}-\bm m_{t,h})\right)\right\|_{\mathrm{op}}
\leq M_{t,h}\quad\text{almost surely},}
\] \CPTSadd{then} \begin{equation}\CPTSadd{\phantomsection\label{eq-remainder-bound}{
\mathbb E(|\widetilde R_{t,h}|\,|\,\mathcal F_t)
\leq M_{t,h}\sum_{i=1}^{p_h}\mathbb E(e_{t+h-i|t}^2|\mathcal F_t).
}}\end{equation}

\end{proposition}

\begin{proof}
\CPTSadd{See Appendix \ref{sec-proof_ar}.}
\end{proof}

It is well-established in the forecasting literature that, for a
zero-mean covariance-stationary time series, the optimal linear
least-squares forecasts have \(h\)-step-ahead errors that are at most
MA\((h-1)\) process \citep{harvey1997}, a property that can be derived
using Wold's representation theorem. \CPTSadd{For the} general nonstationary
autoregressive \CPTSadd{setting considered here}, Proposition~\ref{prp-ma} \CPTSadd{gives
an innovation representation with an explicit remainder}.

\begin{proposition}[Innovation representation of optimal multi-step
forecast errors]\protect\hypertarget{prp-ma}{}\label{prp-ma}

\CPTSadd{Under} the \CPTSadd{assumptions} of Proposition~\ref{prp-ar}\CPTSadd{, for each \(h\in[H]\),}
\begin{equation}\CPTSadd{\phantomsection\label{eq-ma}{
e_{t+h|t}=\varepsilon_{t+h}+
\sum_{i=1}^{h-1}\theta_{t,h,i}\varepsilon_{t+h-i}+U_{t,h},
}}\end{equation} \CPTSadd{where the coefficients \(\theta_{t,h,i}\) are
\(\mathcal F_t\)-measurable and the propagated remainders satisfy
\(U_{t,1}=0\) and
\(U_{t,j}=\widetilde R_{t,j}+\sum_{i=1}^{p_j}\phi_{t,j,i}U_{t,j-i}\) for
\(j\geq2\)}.

\CPTSadd{Moreover, for \(h\geq2\), suppose that the Hessian bounds of
Proposition~\ref{prp-ar} hold and that, for finite deterministic
constants \(M,L,V\geq0\) and all \(2\leq j\leq h\),} \[
\CPTSadd{M_{t,j}\leq M,\qquad
\sum_{i=1}^{p_j}|\phi_{t,j,i}|\leq L,
\qquad \sum_{i=1}^{p_j}\mathbb E(e_{t+j-i|t}^2|\mathcal F_t)\leq V\quad\text{almost surely}.}
\] \CPTSadd{Then} \begin{equation}\CPTSadd{\phantomsection\label{eq-propagated-bound}{
\mathbb E\!\left(|U_{t,h}|\,|\,\mathcal F_t\right)
\leq MV\sum_{k=0}^{h-2}L^k\quad\text{almost surely}.
}}\end{equation}

\end{proposition}

\begin{proof}
\CPTSadd{See Appendix \ref{sec-proof_ma}.}
\end{proof}

\CPTSadd{Dropping the remainders in Propositions \ref{prp-ar} and \ref{prp-ma}
gives linear approximations that motivate regression and moving-average
specifications for multi-step errors, respectively. If the}
autoregressive \CPTSadd{functions are linear} in \CPTSadd{their lagged responses, allowing
intercepts, both remainders vanish and the representations are exact.
Otherwise,} Equation \CPTSadd{\eqref{eq-remainder-bound} controls the Taylor
remainder through curvature and shorter-horizon mean squared errors}, and
\CPTSadd{Equation \eqref{eq-propagated-bound} controls its propagation across
horizons.}

\CPTSadd{Appendix \ref{sec-approximation-errors} shows} that\CPTSadd{, with \(t\), \(h\),
and \(d\) fixed, both remainders have mean absolute size \(O(\delta^2)\)
and} are \CPTSadd{therefore negligible relative to} the \CPTSadd{innovation scale \(\delta\)
as \(\delta\downarrow0\)}. \CPTSadd{This requires deterministic bounds on
gradients and Hessians along} the \CPTSadd{expansion segments that do not depend
on \(\delta\), and innovation conditional variances that are uniformly
\(O(\delta^2)\). The analysis motivates the AcMCP scorecaster} in
\CPTSadd{Section}~\CPTSadd{\ref{sec-novel}}, \CPTSadd{a simplified specification that accounts for
dependence among} forecast errors \CPTSadd{rather than reproducing their} optimal
\CPTSadd{dynamics exactly}.

\subsection{The AcMCP method}\label{sec-novel}

\CPTSadd{Motivated by} these properties, we \CPTSadd{propose} \textbf{autocorrelated
multi-step conformal prediction} (AcMCP). \CPTSadd{Whereas multi-step} extensions
of existing conformal methods \CPTSadd{in} Section~\CPTSadd{\ref{sec-mpid} treat horizons
separately}, AcMCP \CPTSadd{exploits} the \CPTSadd{dependence among} multi-step forecast
errors. \CPTSadd{For each \(h\in[H]\)}, the \(h\)-step-ahead \CPTSadd{threshold} is \CPTSadd{updated
online as} \begin{equation}\CPTSadd{\phantomsection\label{eq-acmcp}{
q_{t+h|t}=q^{\mathrm P}_{t+h|t}+r_t\Bigg(\sum_{i=h+1}^t(\mathrm{err}_{i|i-h}-\alpha)\Bigg)+\tilde{s}_{t+h|t},
}}\end{equation} \CPTSadd{where} the \CPTSadd{lower- and upper-tail thresholds are updated
separately at levels \(\alpha^{-}\) and \(\alpha^{+}\) and combined into}
a \CPTSadd{prediction interval via Equation \eqref{eq-two-sided-pi}. AcMCP
extends} the PID method \CPTSadd{of} \citet{angelopoulos2024} \CPTSadd{in two respects}.

First, we are no longer confined to predicting just one step ahead.
Instead, we can \CPTSadd{produce} multi-step forecasts with accompanying
theoretical coverage guarantees, constructing distribution-free
prediction intervals for \CPTSadd{\(y_{t+1},\dots,y_{t+H}\) from} information
\CPTSadd{available at} time \(t\)

\CPTSadd{Second}, AcMCP\CPTSadd{'s scorecaster accounts for error dependence. Let
\(\tilde e_{t+h|t}\) denote the} forecast of \CPTSadd{\(e_{t+h|t}\), obtained by
equally combining} two simple models: \CPTSadd{(i)} an MA\((h-1)\) model \CPTSadd{fitted to}
the \(h\)-step-ahead errors available up to and including time \(t\),
\CPTSadd{namely \(e_{1+h|1},\dots,e_{t|t-h}\), which captures serial dependence
within a horizon;} and \CPTSadd{(ii)} a linear regression \CPTSadd{of each \(h\)-step-ahead
error} on \CPTSadd{the shorter-horizon} errors from \CPTSadd{the same forecast origin}, which
\CPTSadd{captures dependence} across horizons \CPTSadd{and takes the} recursively \CPTSadd{predicted
errors \(\tilde e_{t+h-1|t},\dots,\tilde e_{t+1|t}\) as inputs}, \CPTSadd{since}
the \CPTSadd{corresponding realized} errors \CPTSadd{are not yet observed at time \(t\)}.
\CPTSadd{The score forecast in Equation \eqref{eq-acmcp} is
\(\tilde s_{t+h|t}=\tilde e_{t+h|t}\) for upper-tail calibration and
\(\tilde s_{t+h|t}=-\tilde e_{t+h|t}\) for lower-tail calibration.}
Importantly, the role of \CPTSadd{\(\tilde e_{t+h|t}\)} is not to forecast the
\CPTSadd{errors precisely}, but to incorporate the autocorrelation structure of
multi-step forecast errors into the resulting prediction intervals.

\CPTSadd{The coverage guarantees below hold under the stated calibration
conditions and do not require the linear approximations in Propositions
\ref{prp-ar} and \ref{prp-ma} to be accurate.}

\subsection{Coverage guarantees}\label{sec-cov}

\CPTSadd{Proposition~\ref{prp-cov_rt} and Corollaries
\ref{cor-cov_qt}, \ref{cor-cov_acmcp} concern single-tail calibration
with fixed target tail miscoverage level \(\alpha\).
Corollary~\ref{cor-cov-two-sided} combines the lower and upper tails at
levels \(\alpha^{-}\) and \(\alpha^{+}\) with
\(\alpha=\alpha^{-}+\alpha^{+}\) the total miscoverage level.}

\begin{proposition}[]\protect\hypertarget{prp-cov_rt}{}\label{prp-cov_rt}

\CPTSadd{Fix \(h\in[H]\) and \(\alpha\in(0,1)\).} Let \CPTSadd{\(\{s_{t+h|t}\}_{t\geq1}\)}
be any sequence of \CPTSadd{finite scores} in \CPTSadd{\((-b,b)\) with \(b\in(0,\infty]\)},
and \CPTSadd{let} \(r_t\) \CPTSadd{satisfy} Equation \eqref{eq-saturation_h} \CPTSadd{and} an
admissible function \(g\). Then the iteration
\(q_{t+h|t}=r_t\left(\sum_{i=h+1}^t(\mathrm{err}_{i|i-h}-\alpha)\right)\)
satisfies\CPTSadd{,} for \CPTSadd{every \(T\geq h+1\)},
\begin{equation}\phantomsection\label{eq-cov_rt}{
\left|\frac{1}{T-h}\sum_{t=h+1}^{T}(\mathrm{err}_{t|t-h}-\alpha)\right| \leq \frac{c \cdot g(T-h) + h}{T-h},
}\end{equation} \CPTSadd{with \(c\)} the constant in Equation
\eqref{eq-saturation_h}. \CPTSadd{Consequently},
\CPTSadd{\(\lim_{T\to\infty}\frac{1}{T-h}\sum_{t=h+1}^{T}\mathrm{err}_{t|t-h}=\alpha\)}.

\end{proposition}

\begin{proof}
See Appendix \ref{sec-proof_cov_rt}.
\end{proof}

\CPTSadd{For} \(h=1\), Proposition~\ref{prp-cov_rt} \CPTSadd{recovers} Proposition 2 of
\citet{angelopoulos2024}. \CPTSadd{With \(n=T-h\) evaluated forecasts}, the bound
\CPTSadd{decomposes into \(c\,g(n)/n\), determined by} the \CPTSadd{integrator's
saturation, and \(h/n\), which captures the feedback delay}. For \CPTSadd{fixed
\(n\)}, \CPTSadd{a longer} horizon \CPTSadd{enlarges} the \CPTSadd{delay term}, \CPTSadd{although actual
calibration errors need not grow. For every fixed \(h\)} the \CPTSadd{bound
vanishes because} \(g\) is sublinear, \CPTSadd{so} the \CPTSadd{delay affects} the
\CPTSadd{finite-sample guarantee but} not \CPTSadd{long-run calibration}.

The \CPTSadd{constant-step} quantile iteration \CPTSadd{is} a \CPTSadd{special case} of
Proposition~\ref{prp-cov_rt}, \CPTSadd{so} its \CPTSadd{calibration bound follows directly}.

\begin{corollary}[]\protect\hypertarget{cor-cov_qt}{}\label{cor-cov_qt}

\CPTSadd{Fix \(h\in[H]\) and \(\alpha\in(0,1)\).} Let \CPTSadd{\(\{s_{t+h|t}\}_{t\geq1}\)}
be any sequence of \CPTSadd{finite scores} in \CPTSadd{\((-b,b)\) with \(0<b<\infty\)}, and
\CPTSadd{let \(q_{2h|h}\in\mathbb R\)} be \CPTSadd{an arbitrary initial threshold}. Then the
iteration \CPTSadd{\(q_{t+h|t}=q_{t+h-1|t-1}+\eta(\mathrm{err}_{t|t-h}-\alpha)\)
with \(\eta>0\)} satisfies\CPTSadd{,} for \CPTSadd{every \(T\geq h+1\)}, \[
\CPTSadd{\left|\frac{1}{T-h}\sum_{t=h+1}^{T}(\mathrm{err}_{t|t-h}-\alpha)\right| \leq \frac{b+|q_{2h|h}|+\eta h}{\eta(T-h)}.}
\] \CPTSadd{In particular,
\(\lim_{T\to\infty}\frac{1}{T-h}\sum_{t=h+1}^{T}\mathrm{err}_{t|t-h}=\alpha\)}.

\end{corollary}

\begin{proof}
See Appendix \ref{sec-proof_cov_qt}.
\end{proof}

\CPTSadd{For} \(h=1\), Corollary~\ref{cor-cov_qt} \CPTSadd{recovers} Proposition 1 of
\citet{angelopoulos2024}. \CPTSadd{The multi-step bound},
\CPTSadd{\((b+|q_{2h|h}|)/(\eta n)+h/n\), separates} the \CPTSadd{effects of score
magnitude, initialization,} and \CPTSadd{learning rate from that of} the \CPTSadd{feedback
delay}. The boundedness \CPTSadd{of} the scores is \CPTSadd{a purely} theoretical \CPTSadd{assumption
that keeps} the \CPTSadd{accumulated} miscoverage error of the \CPTSadd{quantile-tracking
recursion under control}. \CPTSadd{In practice}, the bound itself \CPTSadd{need} not \CPTSadd{be
known}, and the update \CPTSadd{runs unchanged} as long as the scores are finite.

More importantly, Proposition~\ref{prp-cov_rt} also \CPTSadd{establishes
one-sided calibration} for AcMCP. We first \CPTSadd{write} Equation
\eqref{eq-acmcp} as \begin{equation}\phantomsection\label{eq-acmcp_1}{
q_{t+h|t}=\hat{q}_{t+h|t}+r_t\left(\sum_{i=h+1}^t (\mathrm{err}_{i|i-h}-\alpha)\right),
}\end{equation} where \(\hat{q}_{t+h|t}\) \CPTSadd{may} be any function of the
past observations \CPTSadd{\(\{(\bm{x}_i,y_i)\}_{1\leq i\leq t}\)} and \CPTSadd{the earlier
thresholds} \(q_{i+h|i}\)\CPTSadd{, \(i\leq t-1\), and} can \CPTSadd{be viewed} as \CPTSadd{a} forecast
of the \CPTSadd{next threshold from the} available \CPTSadd{history}. \CPTSadd{Taking
\(\hat q_{t+h|t}=q^{\mathrm P}_{t+h|t}+\tilde s_{t+h|t}\) recovers}
AcMCP\CPTSadd{, and taking
\(\hat q_{t+h|t}=q^{\mathrm P}_{t+h|t}+\hat s_{t+h|t}\) recovers MPID}.

\begin{corollary}[]\protect\hypertarget{cor-cov_acmcp}{}\label{cor-cov_acmcp}

\CPTSadd{Fix \(h\in[H]\) and \(\alpha\in(0,1)\).} Let \CPTSadd{the scores \(s_{t+h|t}\)} and
\CPTSadd{baselines \(\hat q_{t+h|t}\)} be \CPTSadd{finite with
\(|s_{t+h|t}-\hat q_{t+h|t}|<b\)} for \CPTSadd{all \(t\geq1\)}, \CPTSadd{where
\(b\in(0,\infty]\),} and \CPTSadd{let} \(r_t\) \CPTSadd{satisfy} Equation
\eqref{eq-saturation_h} \CPTSadd{and} an admissible function \(g\). Then the
\CPTSadd{iteration in Equation \eqref{eq-acmcp_1} satisfies, for every
\(T\geq h+1\),} \[
\CPTSadd{\left|\frac{1}{T-h}\sum_{t=h+1}^{T}(\mathrm{err}_{t|t-h}-\alpha)\right|
\leq \frac{c \cdot g(T-h)+h}{T-h}.}
\] \CPTSadd{In particular,
\(\lim_{T\to\infty}\frac{1}{T-h}\sum_{t=h+1}^{T}\mathrm{err}_{t|t-h}=\alpha\).}

\end{corollary}

\begin{proof}
\CPTSadd{See Appendix \ref{sec-proof_cov_acmcp}.}
\end{proof}

\CPTSadd{Under the stated conditions, integral feedback ensures calibration for
any baseline \(\hat q_{t+h|t}\), and thus without requiring accurate
forecasts from the scorecaster. The scorecaster instead aims to improve
interval efficiency by accounting for error dependence, although this
improvement is not guaranteed by the corollary. Unlike
Corollary~\ref{cor-cov_qt}, which bounds the scores themselves,
Corollary~\ref{cor-cov_acmcp} bounds only their deviations from the
baselines, \(s_{t+h|t}-\hat q_{t+h|t}\).}

\CPTSadd{The next corollary combines the lower- and upper-tail calibration
guarantees to establish coverage for the asymmetric} prediction \CPTSadd{interval
in Equation \eqref{eq-two-sided-pi}.}

\begin{corollary}[]\protect\hypertarget{cor-cov-two-sided}{}\label{cor-cov-two-sided}

\CPTSadd{Fix \(h\in[H]\) and \(\alpha\in(0,1)\). Construct the interval in
Equation \eqref{eq-two-sided-pi} from thresholds} obtained by \CPTSadd{applying}
Equation \eqref{eq-acmcp_1} \CPTSadd{separately to} the \CPTSadd{lower- and upper-tail
scores}, \CPTSadd{with targets \(\alpha^{-},\alpha^{+}>0\) satisfying
\(\alpha^{-}+\alpha^{+}=\alpha\)}. \CPTSadd{Suppose each tail satisfies the
assumptions of Corollary~\ref{cor-cov_acmcp}. Then, for every
\(T\geq h+1\),}
\begin{equation}\CPTSadd{\phantomsection\label{eq-two-sided-finite}{
\frac{1}{T-h}\sum_{t=h+1}^{T}\mathbbm{1}\{y_t\in\hat{\mathcal C}_{t|t-h}^{\alpha}\}
\geq 1-\alpha-2\frac{c\cdot g(T-h)+h}{T-h},
}}\end{equation} \CPTSadd{and hence
\(\liminf_{T\to\infty}\frac{1}{T-h}\sum_{t=h+1}^{T}\mathbbm{1}\{y_t\in\hat{\mathcal C}_{t|t-h}^{\alpha}\}\geq1-\alpha\).}

\CPTSadd{Furthermore, let
\(\chi_{T,h}:=\frac{1}{T-h}\sum_{t=h+1}^{T}\mathbbm{1}\{\hat y_{t|t-h}-q_{t|t-h}^{-}>\hat y_{t|t-h}+q_{t|t-h}^{+}\}\)
be the frequency of crossed endpoints, that is, of empty intervals. Then}
\begin{equation}\CPTSadd{\phantomsection\label{eq-two-sided-enhanced}{
\left|\frac{1}{T-h}\sum_{t=h+1}^{T}\mathbbm{1}\{y_t\in\hat{\mathcal C}_{t|t-h}^{\alpha}\}
-(1-\alpha)\right|\leq 2\frac{c \cdot g(T-h)+h}{T-h}+\chi_{T,h},
}}\end{equation} \CPTSadd{so the empirical coverage converges to \(1-\alpha\)
whenever \(\chi_{T,h}\to0\).}

\end{corollary}

\begin{proof}
See Appendix \CPTSadd{\ref{sec-proof_cov_two_sided}}.
\end{proof}

\CPTSadd{Unlike the preceding single-tail results,
Corollary~\ref{cor-cov-two-sided} provides a lower bound on empirical
coverage of intervals and ensures a limit inferior of at least
\(1-\alpha\), while convergence to exactly \(1-\alpha\) follows when the
endpoint-crossing frequency vanishes. Setting \(\hat q_{t+h|t}=0\), the
result also covers the integral-only iteration in
Proposition~\ref{prp-cov_rt}. Analogous two-sided guarantees for the
constant-step quantile-tracking update shown in
Corollary~\ref{cor-cov_qt} are given in
Corollary~\ref{cor-cov-two-sided-qt} in Appendix
\ref{sec-two-sided-constant-step}.}

\CPTSadd{These results bound empirical coverage at each forecast horizon, both in
finite samples and in the long run, but do not ensure conditional
coverage at individual forecast origins or simultaneous coverage across
horizons. Nor do they ensure local coverage. Appendix
\ref{sec-local-coverage} shows that the coverage shortfall over a window
is bounded by twice the sum of the cumulative calibration bounds at its
endpoints, divided by the window length. This bound vanishes when the
window length grows faster than that sum, but may remain uninformative
for fixed-length rolling windows. Thus, long-run calibration alone does
not guarantee protection against local undercoverage.}

\section{Experiments}\label{sec-experiments}

We evaluate \CPTSadd{MPID and AcMCP} using two simulated settings and two
real-world data sets. \CPTSadd{For benchmarks, we use the multi-step extensions
of split conformal prediction (MSCP), weighted conformal prediction
(MWCP), and adaptive conformal prediction (MACP) described in Appendix
\ref{sec-extensions}, together with the multi-step sequential predictive
conformal inference (SPCI) proposed by \citet{xu2023}.} The following
parameter choices are used:

\begin{enumerate}
\def\labelenumi{(\alph{enumi})}
\tightlist
\item
  we target coverage \(1-\alpha=0.9\)\CPTSadd{, with equal tail miscoverage
  levels \(\alpha^{-}=\alpha^{+}=0.05\)};
\item
  for MWCP, we use \(b=0.99\) as per \citet{barber2023};
\item
  following \citet{angelopoulos2024}, MACP uses a step size
  \(\gamma=0.005\); MPID uses a Theta model as the scorecaster; and both
  MPID and AcMCP use a learning rate of \CPTSadd{\(\eta=0.1\hat{B}_t\)} for
  quantile tracking, where
  \(\hat{B}_t=\max\{s_{t-\Delta+1|t-\Delta-h+1}, \dots, s_{t|t-h}\}\) is
  the highest score over a t\CPTSadd{r}ailing window of length \(\Delta\), which
  is set equal to the calibration set size;
\item
  we adopt a nonlinear saturation function
  \CPTSadd{\(r_t(x)=K_{\mathrm I} \tan \left(x \log (t) / (t C_{\text {sat }})\right)\)},
  where \(\tan (x)=\operatorname{sign}(x) \cdot \infty\) for
  \(x \notin[-\pi / 2, \pi / 2]\), and constants
  \(C_{\text {sat }}, K_{\mathrm{I}}>0\) are chosen heuristically as in
  \citet{angelopoulos2024};
\item
  we use a clipped version of MACP, replacing infinite intervals with
  the largest score observed to date\CPTSadd{;}
\item
  \CPTSadd{for SPCI, we follow the implementation of \citet{xu2023}, using a
  conditional quantile forest with \(10\) trees and maximum depth \(2\)}.
  \CPTSadd{For each horizon, we refit it over the rolling calibration window
  using the latest \(10\) available signed errors as predictors and
  construct} intervals \CPTSadd{at fixed quantile levels \(0.05\)} and \CPTSadd{\(0.95\)}.
  \CPTSadd{All methods share point forecasts and out-of-sample calibration
  errors.}
\end{enumerate}

\subsection{Simulated linear autoregressive process}\label{sec-simar}

We first consider a simulated stationary time series generated from a
simple AR\((2)\) process \[
y_t = 0.8y_{t-1} - 0.5y_{t-2} + \varepsilon_t,
\] where \(\varepsilon_t\) is white noise with error variance
\(\sigma^2 = 1\). After an appropriate burn-in period, we generate
\(N=5000\) data points. Under the sequential split and online learning
settings, we create training sets \(\mathcal{D}_{\text{tr}}\) and
calibration sets \(\mathcal{D}_{\text{cal}}\), each with a length of
\(500\). We use AR\((2)\) models to generate \(1\)- to \(3\)-step-ahead
point forecasts (i.e.~\(H=3\)), using the \texttt{Arima()} function from
the \texttt{forecast} R package \CPTSadd{\citep{hyndman2008, hyndman2024}}. The
goal is to generate prediction intervals using various proposed
conformal prediction methods and evaluate whether they can achieve the
nominal long-run coverage for each separate forecast horizon.

\begin{figure}

\centering{

\pandocbounded{\includegraphics[keepaspectratio]{cpts_files/figure-pdf/fig-AR2_cov-1.pdf}}

}

\caption{\label{fig-AR2_cov}AR(2) simulation results showing rolling
coverage \CPTSadd{(top)} and \CPTSadd{mean} interval width \CPTSadd{(bottom)} for each forecast
horizon. \CPTSadd{Both metrics} are \CPTSadd{computed} over rolling \CPTSadd{windows} of size \(500\).
The \CPTSadd{black dashed horizontal line indicates the} target coverage level
\(1-\alpha=0.9\).}

\end{figure}%

Figure~\ref{fig-AR2_cov} presents the rolling coverage and \CPTSadd{mean} interval
width of each method for each forecast horizon, with metrics computed
over a rolling window of size \(500\). \CPTSadd{For each method and forecast
horizon, rolling median interval widths follow trajectories similar to
those of the corresponding means and are omitted to avoid redundancy.}
The analytic intervals obtained from the \CPTSadd{fitted} AR model are also shown.
We observe that MPID and AcMCP achieve approximately the desired
\(90\%\) coverage level over the rolling windows, while other methods,
including the AR model, undergo much wider swings away from the desired
level, showing high coverage volatility over time. AcMCP constructs
narrower prediction intervals than MPID, despite both methods achieving
similar coverage. Moreover, we see that AcMCP tends to offer
coverage-adaptive prediction intervals, automatically adjusting their
width based on past coverage performance, and results in wider intervals
especially when competing methods undercover, which is to be expected.
In short, AcMCP intervals offer greater adaptivity and more precise
coverage compared to AR, MSCP, MWCP\CPTSadd{, SPCI} and MACP. Both MPID and AcMCP
achieve tight coverage, but AcMCP outputs more informative and smaller
intervals. The benefits of AcMCP are more noticeable when \(h\) grows,
and it does not deteriorate for smaller \(h\). The results can be
further elucidated with Figure~\ref{fig-AR2_box}, which presents
boxplots of rolling coverage and interval width for each method and each
forecast horizon.

\begin{figure}

\centering{

\pandocbounded{\includegraphics[keepaspectratio]{cpts_files/figure-pdf/fig-AR2_box-1.pdf}}

}

\caption{\label{fig-AR2_box}AR(2) simulation results showing boxplots of
the rolling coverage and interval width for each method across different
forecast horizons. \CPTSadd{Black} dashed lines \CPTSadd{indicate} the target coverage level
\CPTSadd{in} the \CPTSadd{left panels and} the median interval width of AcMCP \CPTSadd{in the right
panels}.}

\end{figure}%

\begin{figure}[!b]

\centering{

\pandocbounded{\includegraphics[keepaspectratio]{cpts_files/figure-pdf/fig-AR2_timeplot-1.pdf}}

}

\caption{\label{fig-AR2_timeplot}AR(2) simulation results showing the
prediction interval bounds for the MACP, MPI, and AcMCP methods over a
truncated period of length 500.}

\end{figure}%

\CPTSadd{To examine the effect} of the \CPTSadd{scorecasting} term \CPTSadd{\(\tilde s_{t+h|t}\), we
compare} AcMCP \CPTSadd{with} the version without this term, \CPTSadd{denoted by} MPI.
\CPTSadd{Figure~\ref{fig-AR2_timeplot} displays} their prediction intervals over a
\CPTSadd{selected} period of \(500\) \CPTSadd{observations}. We observe that AcMCP and MPI
indeed construct similar prediction intervals so their lower and upper
bounds mostly overlap with each other. The main differences occur around
time 1320 and during 1470-1500, where AcMCP \CPTSadd{exhibits} a \CPTSadd{localized}
fanning-out effect \CPTSadd{relative} to MPI\CPTSadd{, with wider intervals most apparent
at \(h=3\)}. Figure~\ref{fig-AR2_timeplot} also presents the prediction
interval bounds given by MACP. Since both MACP and AcMCP rely on past
miscoverage information to update prediction intervals, we include MACP
to support a direct comparison of their resulting prediction intervals
while preserving a clear and concise presentation. The prediction
intervals of both AcMCP and MACP can capture certain patterns in the
actual observations, and there is no consistent pattern indicating
dominance of one method over the other in terms of interval width.

\subsection{Simulated nonlinear autoregressive
process}\label{simulated-nonlinear-autoregressive-process}

Next\CPTSadd{,} consider the case of a nonlinear \CPTSadd{data-generating} process, defined
as \[
y_t = \sin(y_{t-1}) + 0.5\log(y_{t-2} + 1) + 0.1y_{t-1}x_{1,t} + 0.3x_{2,t} + \varepsilon_{t},
\] where \(x_{1,t}\) and \(x_{2,t}\) are uniformly distributed on
\([0,1]\), and \(\varepsilon_{t}\) is white noise with variance
\CPTSadd{\(\sigma^2=0.01\)}. Thus, \(y_t\) nonlinearly depends on its lagged
values \(y_{t-1}\), \(y_{t-2}\), and exogenous variables \(x_{1,t}\) and
\(x_{2,t}\).

\begin{figure}[!b]

\centering{

\pandocbounded{\includegraphics[keepaspectratio]{cpts_files/figure-pdf/fig-NL_cov-1.pdf}}

}

\caption{\label{fig-NL_cov}Nonlinear simulation results showing rolling
coverage \CPTSadd{(top)} and \CPTSadd{mean} interval width \CPTSadd{(bottom)} for each forecast
horizon. \CPTSadd{Both metrics} are \CPTSadd{computed} over rolling \CPTSadd{windows} of size \(100\).
The \CPTSadd{black dashed horizontal line indicates the} target coverage level
\(1-\alpha=0.9\).}

\end{figure}%

\begin{figure}[!b]

\centering{

\pandocbounded{\includegraphics[keepaspectratio]{cpts_files/figure-pdf/fig-NL_box-1.pdf}}

}

\caption{\label{fig-NL_box}Nonlinear simulation results showing boxplots
of the rolling coverage and interval width for each method across
different forecast horizons. \CPTSadd{Black} dashed lines \CPTSadd{indicate} the target
coverage level \CPTSadd{in} the \CPTSadd{left panels and} the median interval width of AcMCP
\CPTSadd{in the right panels}.}

\end{figure}%

After an appropriate burn-in period, we generate \(N=2000\) data points.
Under the sequential split and online learning settings, both \CPTSadd{the}
training set \(\mathcal{D}_{\text{tr}}\) and \CPTSadd{the} calibration set
\(\mathcal{D}_{\text{cal}}\) have size \(500\). Given the nonlinear DGP,
we use feed-forward neural networks with a single hidden layer and
lagged inputs to generate \(1\)- to \(3\)-step-ahead point forecasts
(\(H=3\)), using the \texttt{nnetar()} function from the
\texttt{forecast} R package \citep{hyndman2024}. Note that analytic
prediction intervals are unavailable for neural networks, \CPTSadd{so} we do not
include them when presenting the results.

Figure~\ref{fig-NL_cov} shows the rolling coverage and \CPTSadd{mean} interval
width of each method, computed on a rolling window of size \(100\). MPID
and AcMCP are able to maintain minor fluctuations around the target
coverage of \(90\%\) over time, whereas MSCP, MWCP, \CPTSadd{SPCI} and MACP
struggle to sustain the target level and display pronounced fluctuations
over time. Notably, MPID and AcMCP demonstrate greater adaptability,
with higher variability in interval widths compared to competing methods
in order to uphold the desired coverage. The rolling mean interval
widths indicate that AcMCP produces generally narrower intervals than
MPID for \(2\)-step-ahead forecasts but wider intervals for
\(3\)-step-ahead forecasts. Furthermore, under AcMCP, the interval width
at \(h=3\) exceeds that at \(h=2\) in \(53.41\%\) of the test cases,
compared with only \(38.08\%\) under MPID. This pattern suggests more
plausible behavior for AcMCP, as forecast uncertainty typically
increases with the forecast horizon, leading to a greater tendency for
wider prediction intervals.

We provide further insights into the performance of these conformal
prediction methods by presenting boxplots of the rolling coverage and
interval width for each method, as depicted in Figure~\ref{fig-NL_box}.
We observe that coverage variability is higher for MSCP, MWCP\CPTSadd{, SPCI} and
MACP than for MPID and AcMCP, while MPID and AcMCP lead to a lower
effective interval size.

\subsection{Electricity demand data}\label{electricity-demand-data}

We apply the conformal prediction methods using data comprising daily
electricity demand (GW), daily maximum temperature (degrees Celsius),
and holiday information for Victoria, Australia, from 2012 to 2014.
Temperatures are taken from the Melbourne Regional Office weather
station. The left panel of Figure~\ref{fig-elec_data} displays the daily
electricity demand and temperatures over the period. The right panel
shows a nonlinear relationship between demand and temperature, with
demand increasing for low temperatures (due to heating) and increasing
for high temperatures (due to cooling); the two clouds of points
correspond to working days and non-working days.

\begin{figure}[!hb]

\centering{

\pandocbounded{\includegraphics[keepaspectratio]{cpts_files/figure-pdf/fig-elec_data-1.pdf}}

}

\caption{\label{fig-elec_data}Daily electricity demand and corresponding
daily maximum temperatures in 2012--2014, Victoria, Australia.}

\end{figure}%

The response variable is \(\texttt{Demand}\), with two covariates:
\(\texttt{Temperature}\) and \(\texttt{Workday}\) (an indicator variable
for working days). Following \citet{hyndman2021}, we fit a dynamic
regression model with a piecewise linear function of temperature (with a
knot at \(18\) degrees) to generate \(1\)- to \(7\)-step-ahead point
forecasts (\(H=7\)), modeling the errors with an ARIMA process to
capture autocorrelation. We use two years of data for training and
\(100\) observations for calibration.

Figure~\ref{fig-elec_cov} and Figure~\ref{fig-elec_box} compare the
rolling coverage and interval width across method\CPTSadd{s}. These computations
are based on a rolling window of size \(100\). The DR method corresponds
to the analytic intervals obtained from the dynamic regression model.
First, DR consistently attains coverage well above the \(90\%\) target,
resulting in much wider intervals than other methods for \(h=1,2,3,4\).
While this overcoverage improves reliability, the resulting intervals
are overly conservative. Second, MSCP, MWCP, and MACP fail to sustain
the target coverage and noticeably undercover after September 2014 for
all horizons\CPTSadd{. SPCI also produces} narrower intervals than \CPTSadd{MPID} and \CPTSadd{AcMCP
but exhibits substantial undercoverage across horizons. These narrower
intervals therefore come at} the \CPTSadd{cost} of \CPTSadd{substantial undercoverage}.
Third, MPI, MPID, and AcMCP offer wider intervals that largely mitigate
the post--September 2014 undercoverage. Among these, MPID shows slightly
poorer coverage than MPI and AcMCP at \(h=3\), despite producing wider
intervals. Finally, while MPI and AcMCP display similar coverage
pattern\CPTSadd{s}, AcMCP is capable of constructing narrower intervals than MPI,
particularly at larger horizons. Relative to MPID, AcMCP also achieves
greater interval width reduction, with advantages that become more
pronounced as \(h\) increases and without sacrificing performance at
shorter horizons.

\begin{figure}[!hbtp]

\centering{

\pandocbounded{\includegraphics[keepaspectratio]{cpts_files/figure-pdf/fig-elec_cov-1.pdf}}

}

\caption{\label{fig-elec_cov}Electricity demand data results showing
rolling coverage \CPTSadd{(top)} and \CPTSadd{mean} interval width \CPTSadd{(bottom)} for each
forecast horizon. \CPTSadd{Both metrics} are \CPTSadd{computed} over rolling \CPTSadd{windows} of size
\(100\). The \CPTSadd{black dashed horizontal line indicates the} target coverage
level \(1-\alpha=0.9\).}

\end{figure}%

\begin{figure}[!hb]

\centering{

\pandocbounded{\includegraphics[keepaspectratio]{cpts_files/figure-pdf/fig-elec_box-1.pdf}}

}

\caption{\label{fig-elec_box}Electricity demand data results showing
boxplots of the rolling coverage and interval width for each method
across different forecast horizons. \CPTSadd{Black} dashed lines \CPTSadd{indicate} the
target coverage level \CPTSadd{in} the \CPTSadd{left panels and} the median interval width
of AcMCP \CPTSadd{in the right panels}.}

\end{figure}%

\CPTSadd{AcMCP achieves near-nominal overall coverage and stable rolling
coverage, with narrower intervals than competing methods that attain
comparable coverage. These properties support its use in short-term
generation planning: combined with supply forecasts, the intervals can
flag days that warrant closer review of capacity margins and support
assessments of candidate generation plans under different demand
scenarios. However, satisfactory overall and rolling coverage does not
imply conditional coverage under extreme temperatures, high demand, or
other critical operating conditions. The theoretical guarantees concern
empirical coverage at each horizon; they ensure neither coverage
conditional on the information available for a given decision nor
simultaneous coverage across horizons. Costly upper-tail exceedances can
therefore occur despite satisfactory average coverage. Before using the
intervals to set reserves, practitioners should assess the frequency and
magnitude of upper-bound exceedances during high-demand periods, and
weigh this exposure against available generation capacity and the cost
of covering any shortfall.}

\subsection{Eating out expenditure
data}\label{eating-out-expenditure-data}

In our final example, we apply the conformal prediction methods to
forecast the eating out expenditure (\CPTSadd{in millions of dollars}) in
Victoria, Australia. The data include monthly expenditure on cafes,
restaurants and takeaway food services from April 1982 to December 2019,
as shown in Figure~\ref{fig-cafe_data}. The data show an overall upward
trend, \CPTSadd{clear} annual seasonal\CPTSadd{ity}, and \CPTSadd{l}evel-dependent variability.

\begin{figure}[!tb]

\centering{

\pandocbounded{\includegraphics[keepaspectratio]{cpts_files/figure-pdf/fig-cafe_data-1.pdf}}

}

\caption{\label{fig-cafe_data}Monthly expenditure on cafes, restaurants
and takeaway food services in Victoria, Australia, from April 1982 to
December 2019.}

\end{figure}%

We consider three models: ARIMA with logarithmic transformation, ETS,
and STL-ETS \citep{hyndman2021}, and then output their simple average as
final point forecasts. STL-ETS applies ETS to the seasonally adjusted
series obtained via STL decomposition. All models are automatically
trained using the \texttt{forecast} R package \citep{hyndman2024}. We
aim to forecast \(12\) months ahead (\(H=12\)) using \(20\) years of
data for training and \(5\) year\CPTSadd{s} for calibration. Because the test set
contains only \(152\) months, we report coverage and interval width
averaged over the full test period rather than rolling windows.

As the forecast horizons considered in this application are relatively
long, we summarize performance in Figure~\ref{fig-cafe_cov} by reporting
the average coverage gap and average interval width across the whole
test set for each method and each forecast horizon. The results first
show that \CPTSadd{SPCI produces narrower intervals but substantially undercovers
at every horizon, especially longer ones, limiting their usefulness for
uncertainty quantification.} MSCP, MWCP and MACP \CPTSadd{achieve empirical
coverage close to the target} for smaller forecast horizon\CPTSadd{s} but fail to
achieve the desired coverage for larger forecast horizons (\(h>5\)).
Second, for \(h \leq 5\), MPI and AcMCP can approximately achieve the
desired coverage and provide mean interval widths \CPTSadd{comparable to those of
MSCP}, \CPTSadd{MWCP and MACP}. Third, the coverage of MSCP, MWCP\CPTSadd{, SPCI} and MACP
declines gradually as the forecast horizon increases, while MPI and
AcMCP maintain coverage within a tighter range, albeit at the cost of
interval efficiency. Lastly, compared to MPI, AcMCP exhibits slightly
less deviation from the desired coverage across most forecast horizons.

\begin{figure}

\centering{

\pandocbounded{\includegraphics[keepaspectratio]{cpts_files/figure-pdf/fig-cafe_cov-1.pdf}}

}

\caption{\label{fig-cafe_cov}Eating out expenditure data results showing
\CPTSadd{empirical} coverage \CPTSadd{minus the \(90\%\) target level (left)} and \CPTSadd{mean}
interval width \CPTSadd{(right), computed} over the entire test set for each
forecast horizon. The black dashed line indicates \CPTSadd{zero coverage gap}.}

\end{figure}%

\subsection{Post-hoc comparison}\label{post-hoc-comparison}

For MACP, the step size \(\gamma\) controls both the calibration
tightness and the responsiveness of the interval widths: smaller values
lead to smoother interval evolution but slower correction of local
coverage errors, while larger values improve responsiveness at the cost
of increased variability. An analogous trade-off arises in AcMCP, where
the learning rate \(\eta\) governs the balance between adaptivity and
stability. This \CPTSadd{motivates a comparison with} MACP tuned for \CPTSadd{comparable
local coverage variability}.

To examine this, we conduct a post-hoc comparison in which \CPTSadd{the
learning-rate rule and tuning constants} for AcMCP \CPTSadd{are} held fixed, while
\(\gamma\) for MACP is selected via grid search over \([0.005, 0.5]\).
Specifically, \(\gamma\) is chosen to match the sample variance of local
coverage achieved by AcMCP over rolling windows of the test set. \CPTSadd{This is
a post-hoc comparison based on the test set, rather than a prospective
tuning procedure or a guarantee of optimality.}

Table~\ref{tbl-posthoc} summarizes, for each forecast horizon, the mean
coverage over the full test set, the minimum and maximum local coverage
across rolling windows, and the mean and median interval widths for MACP
and AcMCP across the data sets considered above. In simple linear and
nonlinear simulations, both methods attain mean coverage close to the
nominal \CPTSadd{\(90\%\)}, with AcMCP consistently producing comparable or
narrower intervals when coverage is comparable. More pronounced
differences arise in real data sets and at longer horizons. For
electricity demand, MACP displays systematic undercoverage as the
horizon increases, particularly in terms of minimum local coverage,
whereas AcMCP maintains more stable coverage around the target by
adaptively adjusting interval widths. A similar pattern appears for
eating out expenditure data, where MACP suffers substantial coverage
degradation at larger horizons despite tuning, while AcMCP preserves
near-nominal coverage, albeit with wider but informative intervals.
Overall, these results indicate that AcMCP provides superior coverage
control and, when coverage is comparable, yields prediction intervals
that are shorter or at least comparable in length.

\begin{table}

\caption{\label{tbl-posthoc}Performance of MACP and AcMCP across the four datasets, summarized by mean test-set coverage, minimum and maximum local coverage over rolling windows, and mean and median interval widths for each forecast horizon.}

\centering{

\resizebox{\linewidth}{!}{
\begin{tabular}{lrrrrrrrrrr}
\toprule
     & \multicolumn{2}{c}{Mean coverage (\%)}               & \multicolumn{2}{c}{Min. local coverage (\%)}              & \multicolumn{2}{c}{Max. local coverage (\%)}              & \multicolumn{2}{c}{Mean interval width}              & \multicolumn{2}{c}{Median interval width}            \\
\cmidrule(l{3pt}r{3pt}){2-3} \cmidrule(l{3pt}r{3pt}){4-5} \cmidrule(l{3pt}r{3pt}){6-7} \cmidrule(l{3pt}r{3pt}){8-9} \cmidrule(l{3pt}r{3pt}){10-11}
     & \multicolumn{1}{c}{MACP} & \multicolumn{1}{c}{AcMCP} & \multicolumn{1}{c}{MACP} & \multicolumn{1}{c}{AcMCP} & \multicolumn{1}{c}{MACP} & \multicolumn{1}{c}{AcMCP} & \multicolumn{1}{c}{MACP} & \multicolumn{1}{c}{AcMCP} & \multicolumn{1}{c}{MACP} & \multicolumn{1}{c}{AcMCP} \\
\midrule
     & \multicolumn{10}{c}{AR(2) simulation data}                                                                                                                                                                                                                                       \\[0.3em]
h=1  & 89.95                    & 90.00                     & 89.20                    & 89.40                     & 90.80                    & 90.80                     & 4.04                     & 3.55                      & 3.62                     & 3.54                      \\
h=2  & 89.89                    & 89.94                     & 89.00                    & 89.00                     & 91.00                    & 91.00                     & 5.08                     & 4.68                      & 4.85                     & 4.66                      \\
h=3  & 89.94                    & 90.02                     & 88.80                    & 88.80                     & 91.00                    & 91.20                     & 5.16                     & 4.68                      & 5.00                     & 4.84                      \\
\cmidrule(l{3pt}r{3pt}){2-11}
     & \multicolumn{10}{c}{Nonlinear simulation data}                                                                                                                                                                                                                                   \\[0.3em]
h=1  & 90.00                    & 90.00                     & 87.00                    & 86.00                     & 94.00                    & 93.00                     & 0.35                     & 0.36                      & 0.33                     & 0.36                      \\
h=2  & 89.98                    & 90.18                     & 86.00                    & 87.00                     & 94.00                    & 93.00                     & 0.37                     & 0.37                      & 0.33                     & 0.36                      \\
h=3  & 89.86                    & 90.06                     & 87.00                    & 87.00                     & 94.00                    & 94.00                     & 0.35                     & 0.38                      & 0.33                     & 0.38                      \\
\cmidrule(l{3pt}r{3pt}){2-11}
     & \multicolumn{10}{c}{Electricity demand data}                                                                                                                                                                                                                                     \\[0.3em]
h=1  & 89.92                    & 89.15                     & 86.00                    & 85.00                     & 91.00                    & 92.00                     & 25.02                    & 24.44                     & 22.52                    & 24.24                     \\
h=2  & 88.67                    & 90.23                     & 83.00                    & 86.00                     & 91.00                    & 94.00                     & 30.56                    & 30.72                     & 29.03                    & 28.87                     \\
h=3  & 88.58                    & 90.94                     & 81.00                    & 86.00                     & 92.00                    & 93.00                     & 34.81                    & 34.35                     & 31.64                    & 33.57                     \\
h=4  & 89.29                    & 90.87                     & 84.00                    & 88.00                     & 93.00                    & 94.00                     & 41.93                    & 40.74                     & 34.90                    & 39.48                     \\
h=5  & 89.60                    & 94.80                     & 86.00                    & 91.00                     & 93.00                    & 99.00                     & 41.17                    & 55.62                     & 34.72                    & 45.43                     \\
h=6  & 90.32                    & 90.73                     & 85.00                    & 88.00                     & 92.00                    & 97.00                     & 36.34                    & 46.65                     & 35.08                    & 41.56                     \\
h=7  & 87.80                    & 89.84                     & 84.00                    & 87.00                     & 90.00                    & 95.00                     & 36.55                    & 47.45                     & 35.02                    & 47.98                     \\
\cmidrule(l{3pt}r{3pt}){2-11}
     & \multicolumn{10}{c}{Eating out expenditure data}                                                                                                                                                                                                                                 \\[0.3em]
h=1  & 85.62                    & 88.89                     & 80.00                    & 83.33                     & 90.00                    & 93.33                     & 76.70                    & 65.51                     & 64.21                    & 63.08                     \\
h=2  & 84.77                    & 90.01                     & 78.33                    & 83.33                     & 93.33                    & 96.67                     & 97.93                    & 91.50                     & 84.26                    & 89.10                     \\
h=3  & 80.54                    & 89.26                     & 68.33                    & 83.33                     & 95.00                    & 95.00                     & 104.10                   & 112.42                    & 101.46                   & 107.83                    \\
h=4  & 59.18                    & 89.12                     & 15.00                    & 81.67                     & 91.67                    & 98.33                     & 89.74                    & 116.80                    & 77.68                    & 111.16                    \\
h=5  & 46.90                    & 91.03                     & 13.33                    & 81.67                     & 65.00                    & 96.67                     & 59.72                    & 126.12                    & 68.74                    & 120.97                    \\
h=6  & 36.36                    & 90.91                     & 8.33                     & 85.00                     & 41.67                    & 96.67                     & 29.26                    & 178.17                    & 62.05                    & 178.60                    \\
h=7  & 53.90                    & 86.52                     & 31.67                    & 73.33                     & 61.67                    & 98.33                     & 63.10                    & 186.58                    & 81.44                    & 198.27                    \\
h=8  & 47.48                    & 88.29                     & 23.33                    & 83.33                     & 58.33                    & 98.33                     & 55.30                    & 192.07                    & 78.62                    & 192.59                    \\
h=9  & 32.12                    & 89.78                     & 1.67                     & 86.67                     & 48.33                    & 100.00                    & 20.29                    & 194.40                    & 47.47                    & 168.96                    \\
h=10 & 71.85                    & 89.63                     & 53.33                    & 85.00                     & 90.00                    & 93.34                    & 116.88                   & 195.83                    & 121.56                   & 193.68                    \\
h=11 & 69.92                    & 88.72                     & 41.67                    & 85.00                     & 98.33                    & 96.67                     & 102.20                   & 233.80                    & 99.26                    & 220.92                    \\
h=12 & 69.47                    & 91.60                     & 65.00                    & 83.33                     & 83.33                    & 98.33                     & 127.19                   & 336.02                    & 128.28                   & 297.48                    \\
\bottomrule
\end{tabular}}

}

\end{table}%

\section{Conclusion and discussion}\label{conclusion-and-discussion}

We \CPTSadd{propose MPID and AcMCP to construct} prediction \CPTSadd{sets for multi-step}
time series forecasting. \CPTSadd{Both} methods \CPTSadd{extend conformal PID control
\citep{angelopoulos2024}} to \CPTSadd{accommodate delayed feedback inherent in
this setting}, \CPTSadd{while AcMCP additionally incorporates dependence among}
forecast errors \CPTSadd{across horizons into} its \CPTSadd{scorecaster}. \CPTSadd{Decompositions of
optimal forecast errors, together with explicit bounds} on \CPTSadd{Taylor
remainders}, \CPTSadd{provide} a \CPTSadd{structural rationale} for \CPTSadd{this design} and \CPTSadd{clarify
the conditions supporting the underlying linear approximations. For both
methods, we establish finite-sample bounds and horizon-specific} long-run
guarantees \CPTSadd{for empirical coverage under the stated calibration
conditions,} without \CPTSadd{requiring exchangeability, stationarity, or}
assumptions on \CPTSadd{how the} data \CPTSadd{distribution evolves over time}. \CPTSadd{These
guarantees do not depend on the accuracy of the linear approximations.
The coverage bounds distinguish the contribution of calibration feedback
from that of forecast delay, whose effect vanishes asymptotically for
each fixed horizon.} Across \CPTSadd{the experiments considered}, \CPTSadd{MPID and} AcMCP
\CPTSadd{adapt their interval widths} to \CPTSadd{changing time series dynamics, generally
maintaining more stable rolling coverage around} the target \CPTSadd{than the
benchmark methods. Among comparably well-calibrated methods, AcMCP often
produces narrower} intervals\CPTSadd{, demonstrating the empirical value of
incorporating forecast error dependence into conformal calibration}.

Our analysis \CPTSadd{of forecast errors} is restricted to an ex-post forecasting
setting, in which future values of exogenous predictors are assumed to
be \CPTSadd{known}. In this setting, interval width reflects both \emph{aleatoric
uncertainty}, arising from intrinsic randomness in \CPTSadd{the} DGP, and
\emph{epistemic uncertainty}, stemming from a lack of knowledge about
the underlying DGP \citep{sale2025}. In many practical forecasting
scenarios, however, predictors must themselves be forecast. \CPTSadd{Uncertainty
in future predictor values} and \CPTSadd{in} the \CPTSadd{models used to forecast them can
then propagate into response} forecasts, \CPTSadd{motivating explicit treatment} of
\CPTSadd{predictor} uncertainty \CPTSadd{in conformal inference.}

\CPTSadd{Although} the \CPTSadd{forecast error analysis assumes known future predictors,}
the \CPTSadd{coverage guarantees remain applicable when these predictors are
forecast}, \CPTSadd{provided that the assumptions of Corollary~\ref{cor-cov_acmcp}}
and \CPTSadd{Corollary~\ref{cor-cov-two-sided} hold. Appendix
\ref{sec-predictor-uncertainty} compares AcMCP using known and forecast
future predictors} in \CPTSadd{a simulation. Both settings achieve near-nominal
empirical} coverage and \CPTSadd{small rolling coverage errors}, \CPTSadd{but forecasting}
the \CPTSadd{predictors produces wider intervals, with larger increases at higher
levels of predictor uncertainty and longer forecast horizons}.

These \CPTSadd{considerations} suggest several directions for future research. One
is to improve interval efficiency while preserving coverage guarantees.
Our methods \CPTSadd{prioritize} efficient online implementation through simple
\CPTSadd{updates}. \CPTSadd{A complementary approach is} Bellman conformal inference
\citep{yang2024ts}, which \CPTSadd{uses a look-ahead} optimization to \CPTSadd{balance}
interval length \CPTSadd{against a} miscoverage \CPTSadd{penalty}, \CPTSadd{but issues only a
one-step-ahead interval}. \CPTSadd{Extending this optimization principle} to
\CPTSadd{multi-step interval construction} while retaining coverage guarantees is
\CPTSadd{a} promising direction\CPTSadd{. Automated learning-rate selection could further
reduce manual tuning and improve adaptation to changing} forecasting
\CPTSadd{conditions. Other avenues include establishing conditional coverage
guarantees under explicit restrictions on distributional change and
developing a theoretical understanding of how} uncertainty in future
predictors \CPTSadd{and their forecasting models affects interval efficiency and
local calibration}.

% \newpage
\appendix
% \pagenumbering{arabic}% resets `page` counter to 1
\setcounter{section}{0}
\renewcommand{\thesection}{Appendix \Alph{section}}
\renewcommand{\thesubsection}{\Alph{section}.\arabic{subsection}}
% \renewcommand{\thefigure}{A\arabic{figure}}
% \renewcommand{\thetable}{A\arabic{table}}
% \setcounter{figure}{0}
% \setcounter{table}{0}

\section*{Acknowledgments}\label{acknowledgments}
\addcontentsline{toc}{section}{Acknowledgments}

This work was supported by the \CPTSadd{National Natural Science Foundation of
China} (No.~\CPTSadd{72601399),} the Presidential Foundation of the Academy of
Mathematics and Systems Science, Chinese Academy of Sciences, China
\CPTSadd{(}No.~E555930101\CPTSadd{), and the Australian Research Council Industrial
Transformation Training Centre in Optimisation Technologies, Integrated
Methodologies, and Applications (OPTIMA) (No}.\CPTSadd{~IC200100009).}

\section*{Disclosure statement}\label{disclosure-statement}
\addcontentsline{toc}{section}{Disclosure statement}

The authors report there are no competing interests to declare.

\section*{Supplementary materials}\label{supplementary-materials}
\addcontentsline{toc}{section}{Supplementary materials}

\CPTSadd{The Appendix provides} proofs of \CPTSadd{all} propositions and corollaries\CPTSadd{,
additional multi-step conformal methods, analyses of approximation
errors and local coverage, supplementary two-sided coverage results, and
a simulation examining} the \CPTSadd{effect of uncertainty} in \CPTSadd{future predictors}.

\section*{References}\label{references}
\addcontentsline{toc}{section}{References}

\begingroup

\renewcommand{\bibsection}{}
\bibliography{references.bib}

\endgroup

\newpage
\appendix
\pagenumbering{arabic}% resets `page` counter to 1

\setcounter{section}{0}
\setcounter{subsection}{0}
\setcounter{figure}{0}
\setcounter{table}{0}

\renewcommand{\thesection}{\Alph{section}}
\renewcommand{\thesubsection}{A\arabic{subsection}}
\renewcommand{\thefigure}{A\arabic{figure}}
\renewcommand{\thetable}{A\arabic{table}}

\section*{Appendix}\label{appendix}
\addcontentsline{toc}{section}{Appendix}

\subsection{\CPTSadd{Multi-step extensions of existing conformal
methods}}\label{sec-extensions}

\subsubsection{\CPTSadd{Online multi-step split conformal
prediction}}\label{online-multi-step-split-conformal-prediction}

\CPTSadd{Split conformal prediction \citep[SCP,][]{papadopoulos2002, vovk2005} is
a holdout method for building prediction sets using a model pre-trained
on a training set. A key advantage of SCP is its ability to guarantee
coverage by assuming data exchangeability. In a regression setting, SCP
randomly separates} the \CPTSadd{available \(n\) data points,
\(Z_i = (X_i, Y_i) \in \mathbb{R}^d \times \mathbb{R}\),
\(i=1, \dots, n\), into a proper training set of size \(n_t\) and a
calibration set of size \(n_c\). Given a regression model
\(\hat{\mu}: \mathbb{R}^d \rightarrow \mathbb{R}\) fitted on the
training set and a score function \(\mathcal{S}\), nonconformity scores
\(s_i = \mathcal{S}\left(X_i, Y_i\right)\) are computed on the
calibration set to measure the nonconformity between the calibration
response values and the predicted values obtained from the fitted model
\(\hat{\mu}\). Then SCP computes the prediction set for the test
response \(Y_{n+1}\) using} \[
\CPTSadd{\hat{\mathcal{C}}_{n+1}^{\alpha}\left(X_{n+1}\right) = \left\{y\in\mathbb{R}: \mathcal{S}\left(X_{n+1}, y\right) \leq Q_{1-\alpha}\left(\sum_{i \in \mathcal{D}_{\text{cal}}}\frac{1}{n_c+1}\cdot\delta_{s_{i}}+\frac{1}{n_c+1}\cdot\delta_{+\infty}\right)\right\},}
\] \CPTSadd{where \(Q_\tau(\cdot)\) denotes the \(\tau\)-quantile of its
argument, and \(\delta_a\) denotes the point mass at \(a\). Time series
data are inherently nonexchangeable due to their temporal dependence and
autocorrelation. Therefore, directly applying SCP to} time series \CPTSadd{data
would violate the method's exchangeability assumption and compromise its
coverage guarantee.}

\CPTSadd{Here we introduce online} \textbf{\CPTSadd{multi-step split conformal prediction}}
\CPTSadd{(MSCP) as} a \CPTSadd{generalization of SCP \citep{papadopoulos2002, vovk2005} to
recursively update all \(h\)-step-ahead prediction intervals over time.
MSCP applies conformal inference in an online fashion, updating
prediction intervals} as \CPTSadd{new data points are received. Specifically, for
each \(h \in [H]\), the one-sided prediction set is}
\begin{equation}\CPTSadd{\phantomsection\label{eq-mscp}{
\hat{\mathcal{C}}_{t+h|t}^{\alpha,\mathrm{tail}} = \Bigg\{y\in\mathbb{R}: s_{t+h|t}^{y} \leq Q_{1-\alpha}\Bigg(\sum_{i=t-t_c+1}^{t}\frac{1}{t_c+1}\cdot\delta_{s_{i|i-h}}+\frac{1}{t_c+1}\cdot\delta_{+\infty}\Bigg)\Bigg\},
}}\end{equation} \CPTSadd{where \(s_{t+h|t}^{y}\) is the relevant signed
nonconformity score} defined in \CPTSadd{Section~\ref{sec-setup}. Applying}
Equation \CPTSadd{\eqref{eq-mscp} to both tails and combining the resulting
thresholds as in Equation \eqref{eq-two-sided-pi} yields the MSCP
interval}.

\subsubsection{\CPTSadd{Online multi-step weighted conformal
prediction}}\label{online-multi-step-weighted-conformal-prediction}

\CPTSadd{\citet{barber2023} propose a nonexchangeable conformal prediction
procedure (NexCP) that generalizes SCP to certain nonexchangeable
settings by assigning higher weights to observations believed to share}
the \CPTSadd{same distribution as the test data. NexCP assumes these weights are
fixed and data-independent. Under exchangeability, NexCP retains the
same coverage guarantees as SCP; when exchangeability is violated, the
coverage gap is governed} by \CPTSadd{the total variation distance between swapped
nonconformity score vectors}, \CPTSadd{which can be substantial in time series
settings.}

\CPTSadd{The online} \textbf{\CPTSadd{multi-step weighted conformal prediction}} \CPTSadd{(MWCP)
method we propose here adapts the NexCP method} to \CPTSadd{the online setting for}
time \CPTSadd{series forecasting. MWCP uses weighted quantile estimates for
constructing prediction intervals}, \CPTSadd{contrasting with the MSCP definitions
where all nonconformity scores for calibration are implicitly assigned
equal weight.}

\CPTSadd{For the subsequent empirical study, we choose fixed weights
\(w_i = b^{t+1-i}\), \(b \in (0, 1)\)} and \CPTSadd{\(i=t-t_c+1, \dots, t\), for
nonconformity scores on} the corresponding \CPTSadd{calibration set. In this
setting, weights decay exponentially as the nonconformity scores get
older, akin to the rationale behind the simple exponential smoothing
method in time series forecasting \citep{hyndman2021}. Then, for each
\(h \in [H]\), MWCP updates the} \(h\)-step-ahead \CPTSadd{one-sided prediction
set:} \[
\CPTSadd{\hat{\mathcal{C}}_{t+h|t}^{\alpha,\mathrm{tail}} = \Bigg\{y\in\mathbb{R}: s_{t+h|t}^{y} \leq Q_{1-\alpha}\Bigg(\sum_{i=t-t_c+1}^{t}\tilde{w}_i\cdot\delta_{s_{i|i-h}}+\tilde{w}_{t+1}\cdot\delta_{+\infty}\Bigg)\Bigg\},}
\] \CPTSadd{where \(\tilde{w}_i\) and \(\tilde{w}_{t+1}\) are normalized weights
given by} \[
\CPTSadd{\tilde{w}_i = \frac{w_i}{\sum_{i=t-t_c+1}^{t}w_i+1}, \text{ for } i \in \{t-t_c+1, \dots, t\} \quad \text{and} \quad \tilde{w}_{t+1} =  \frac{1}{\sum_{i=t-t_c+1}^{t}w_i+1}.}
\] \CPTSadd{As with MSCP, the two one-sided MWCP sets are combined according to
Equation \eqref{eq-two-sided-pi}}.

\subsubsection{\CPTSadd{Multi-step adaptive conformal
prediction}}\label{multi-step-adaptive-conformal-prediction}

\CPTSadd{Next we extend the adaptive conformal inference (ACI) method proposed by
\citet{gibbs2021} to address multi-step time series forecasting,
introducing the} \textbf{\CPTSadd{multi-step adaptive conformal prediction}} \CPTSadd{(MACP)
method}. \CPTSadd{For a generic tail with fixed target miscoverage level
\(\alpha\), the method adapts the nominal level \(\alpha_{t+h|t}\) used
for quantile estimation to past tail miscoverage indicators.
Specifically, for each \(h \in [H]\),} we \CPTSadd{update \(\alpha_{t+h|t}\)
according to} \begin{equation}\phantomsection\label{eq-macp}{
\alpha_{t+h|t} := \alpha_{t+h-1|t-1} + \gamma(\alpha - \mathrm{err}_{t|t-h}).
}\end{equation} \CPTSadd{The corresponding one-sided set is computed from
Equation \eqref{eq-mscp} with \(\alpha\) replaced by \(\alpha_{t+h|t}\).}
For \CPTSadd{levels at or below zero}, \CPTSadd{the one-sided set is \(\mathbb R\); for
levels at or above one}, \CPTSadd{it} is \CPTSadd{empty}. \CPTSadd{Here, \(\gamma > 0\) is a fixed
step size, and \(\alpha_{2h|h}\) is} the \CPTSadd{initial estimate, typically set
to \(\alpha\). The upper- and lower-tail recursions} are \CPTSadd{run separately
and then combined through Equation \eqref{eq-two-sided-pi}}.

\CPTSadd{Equation \eqref{eq-macp} uses the tail miscoverage indicator observed at
time \(t\) to update the forecast for time \(t+h\). The nominal tail
miscoverage level increases when the observation satisfies the relevant
upper or lower bound (\(s_{t|t-h}\leq q_{t|t-h}\)) and decreases when it
violates that bound (\(s_{t|t-h}>q_{t|t-h}\)).} For \CPTSadd{a fixed calibration
sample}, \CPTSadd{decreasing this level widens} the \CPTSadd{one-sided prediction set.
Equation \eqref{eq-macp} reduces to the ACI update when \(h=1\).}

\CPTSadd{We do not use the update
\(\alpha_{t+1|t-h+1} := \alpha_{t|t-h} + \gamma(\alpha - \mathrm{err}_{t|t-h})\)},
\CPTSadd{because it requires \(y_t\)}, \CPTSadd{which is unavailable at} the \CPTSadd{forecast origin
\(t-h+1\) when \(h>1\).}

\CPTSadd{Selecting the parameter \(\gamma\) is pivotal yet challenging.
\citet{gibbs2021} relate the choice} of \CPTSadd{\(\gamma\) to the rate} of
\CPTSadd{distributional change. Several strategies have been proposed to avoid}
the \CPTSadd{necessity of selecting \(\gamma\). For example, \citet{zaffran2022}
use an adaptive aggregation of multiple ACIs} with \CPTSadd{a set of candidate
values for \(\gamma\), determining weights based on their historical
performance. \citet{bastani2022} propose a multivalid prediction
algorithm in which the prediction set is established by selecting a
threshold from a sequence of candidate thresholds. Both methods rely on
update schemes that place substantial weight on older historical data,
which may impede rapid adaptation} to \CPTSadd{abrupt changes \citep{gibbs2024}.
Therefore, \citet{gibbs2024} propose an alternative expert selection
scheme that adaptively tunes} the \CPTSadd{step-size parameter over time and
places greater emphasis on more recent observations by construction,
enabling faster responses to sudden environmental shifts.}

\CPTSadd{The theoretical coverage properties of ACI suggest that a larger value
for \(\gamma\) generally results} in \CPTSadd{less deviation from the target
coverage}. \CPTSadd{As there is no restriction on \(\alpha_{t+h|t}\), it can drift
below \(0\) or above \(1\); a larger \(\gamma\) may therefore lead to
frequent empty or unbounded one-sided sets as the procedure responds to
whether the relevant bound is violated.}

\subsection{\texorpdfstring{Proof of
Proposition~\ref{prp-ar}}{Proof of Proposition~}}\label{sec-proof_ar}

\begin{proof}
\CPTSadd{Fix a forecast origin \(t\). For \(h=1\)}, \CPTSadd{all lagged responses and
required predictors are \(\mathcal F_t\)-measurable. Since
\(\mathbb E(\varepsilon_{t+1}|\mathcal F_t)=0\)}, we \CPTSadd{have} \[
\CPTSadd{\hat y_{t+1|t}=f_{t+1}(y_{(t+1-d):t},\bm x_{(t+1-k):(t+1)})}
\] \CPTSadd{and hence \(e_{t+1|t}=\varepsilon_{t+1}\), with
\(\widetilde R_{t,1}=0\).}

\CPTSadd{For \(h\geq2\),} write
\CPTSadd{\(\bm\Delta_{t,h}=\bm z_{t,h}-\bm m_{t,h}=(e_{t+h-1|t},\dots,e_{t+h-p_h|t})^{\prime}\).
Applying a first-order Taylor expansion and retaining the remainder
gives} \begin{align*}
\CPTSadd{y_{t+h}}
&\CPTSadd{=F_{t,h}(\bm z_{t,h})+\varepsilon_{t+h}}\\
&\CPTSadd{=F_{t,h}(\bm m_{t,h})
+\sum_{i=1}^{p_h}\partial_iF_{t,h}(\bm m_{t,h})e_{t+h-i|t}
+R_{t,h}+\varepsilon_{t+h},}
\end{align*} \CPTSadd{where \(R_{t,h}\) is the Taylor remainder beyond first
order.}

\CPTSadd{The expansion point and derivative coefficients are
\(\mathcal F_t\)-measurable. Since
\(\mathbb E(e_{t+h-i|t}|\mathcal F_t)=0\) for \(1\leq i\leq p_h\) and},
\CPTSadd{by} the \CPTSadd{tower property, \(\mathbb E(\varepsilon_{t+h}|\mathcal F_t)=0\),
taking conditional expectations under the stated integrability
assumptions gives} \[
\CPTSadd{\hat y_{t+h|t}=F_{t,h}(\bm m_{t,h})+\mathbb E(R_{t,h}|\mathcal F_t).}
\] \CPTSadd{Subtracting this identity from the expansion yields} \[
\CPTSadd{e_{t+h|t}=\varepsilon_{t+h}+\sum_{i=1}^{p_h}\phi_{t,h,i}e_{t+h-i|t}+\widetilde R_{t,h},}
\] \CPTSadd{where \(\widetilde R_{t,h}=R_{t,h}-\mathbb E(R_{t,h}|\mathcal F_t)\)
and \(\mathbb E(\widetilde R_{t,h}|\mathcal F_t)=0\). Together with the
case \(h=1\), interpreting} an \CPTSadd{empty sum as zero}, \CPTSadd{this proves Equation
\eqref{eq-ar}. If \(f_{t+h}\) is linear in its lagged responses,
allowing an intercept, then \(R_{t,h}=\widetilde R_{t,h}=0\).}

\CPTSadd{To establish Equation \eqref{eq-remainder-bound}, use the integral
representation} \[
\CPTSadd{R_{t,h}=\int_0^1(1-u)\bm\Delta_{t,h}^{\prime}
\nabla^2 F_{t,h}(\bm m_{t,h}+u\bm\Delta_{t,h})\bm\Delta_{t,h}\,du,}
\] where \CPTSadd{\(\nabla^2F_{t,h}\) is the Hessian with respect to the
unobserved response coordinates}. \CPTSadd{The Hessian bound} implies
\CPTSadd{\(|R_{t,h}|\leq(M_{t,h}/2)\|\bm\Delta_{t,h}\|_2^2\). Conditional
Jensen's inequality and} the \CPTSadd{triangle inequality give} \[
\CPTSadd{\mathbb E\!\left[|\widetilde R_{t,h}|\,|\,\mathcal F_t\right]
\leq2\mathbb E\!\left[|R_{t,h}|\,|\,\mathcal F_t\right]
\leq M_{t,h}\mathbb E\!\left[\|\bm\Delta_{t,h}\|_2^2\,|\,\mathcal F_t\right].}
\] \CPTSadd{Since \(\|\bm\Delta_{t,h}\|_2^2=\sum_{i=1}^{p_h}e_{t+h-i|t}^2\), this
proves Equation \eqref{eq-remainder-bound}}.
\end{proof}

\subsection{\texorpdfstring{Proof of
Proposition~\ref{prp-ma}}{Proof of Proposition~}}\label{sec-proof_ma}

\begin{proof}
\CPTSadd{Fix a forecast origin \(t\). By} Proposition~\ref{prp-ar}\CPTSadd{,} \begin{align*}
\CPTSadd{e_{t+1|t}}&\CPTSadd{=\varepsilon_{t+1},}\\
\CPTSadd{e_{t+j|t}}&\CPTSadd{=\varepsilon_{t+j}+\sum_{i=1}^{p_j}\phi_{t,j,i}e_{t+j-i|t}+\widetilde R_{t,j},\qquad 2\leq j\leq h.}
\end{align*}

\CPTSadd{Recursively substituting the earlier-horizon relations gives} \[
\CPTSadd{e_{t+h|t}=\varepsilon_{t+h}+\sum_{i=1}^{h-1}\theta_{t,h,i}\varepsilon_{t+h-i}+U_{t,h}.}
\] \CPTSadd{The coefficients \(\theta_{t,h,i}\) are finite sums of products of
the \(\phi\) coefficients and are therefore \(\mathcal F_t\)-measurable}.
\CPTSadd{Collecting the remainder terms yields}
\begin{equation}\CPTSadd{\phantomsection\label{eq-propagated-remainder}{
U_{t,1}=0,\qquad
U_{t,j}=\widetilde R_{t,j}+\sum_{i=1}^{p_j}\phi_{t,j,i}U_{t,j-i},\qquad 2\leq j\leq h.
}}\end{equation} \CPTSadd{This proves Equation \eqref{eq-ma} without discarding
any Taylor remainder.}

\CPTSadd{Under the additional assumptions of} Proposition~\CPTSadd{\ref{prp-ma}, Equation
\eqref{eq-remainder-bound} gives
\(\mathbb E(|\widetilde R_{t,j}|\,|\,\mathcal F_t)\leq MV\) for
\(2\leq j\leq h\). The triangle inequality} and
\CPTSadd{\(\mathcal F_t\)-measurability of the coefficients imply} \[
\CPTSadd{\mathbb E(|U_{t,j}|\,|\,\mathcal F_t)
\leq MV+\sum_{i=1}^{p_j}|\phi_{t,j,i}|
\mathbb E(|U_{t,j-i}|\,|\,\mathcal F_t).}
\] \CPTSadd{Let \(G_1=0\) and \(G_j=\sum_{k=0}^{j-2}L^k\) for \(j\geq2\). The
bound \(\mathbb E(|U_{t,1}|\,|\,\mathcal F_t)\leq MVG_1\) holds since
\(U_{t,1}=0\). If it holds for all horizons below \(j\)}, \CPTSadd{then}
\begin{align*}
\CPTSadd{\mathbb E(|U_{t,j}|\,|\,\mathcal F_t)}
&\CPTSadd{\leq MV\left(1+\sum_{i=1}^{p_j}|\phi_{t,j,i}|G_{j-i}\right)}\\
&\CPTSadd{\leq MV(1+LG_{j-1})=MVG_j,}
\end{align*} where \CPTSadd{\(G_{j-i}\leq G_{j-1}\) because \(L\geq0\). Strong
induction therefore establishes the bound for all \(1\leq j\leq h\), and
taking \(j=h\) proves Equation \eqref{eq-propagated-bound}.}
\end{proof}

\subsection{\texorpdfstring{Approximation error analysis for
Propositions \ref{prp-ar} and
\ref{prp-ma}}{Approximation error analysis for Propositions  and }}\label{sec-approximation-errors}

\CPTSadd{Omitting the remainders in Propositions \ref{prp-ar} and \ref{prp-ma}
yields linear approximations with errors \(\widetilde R_{t,h}\) and
\(U_{t,h}\). Equation \eqref{eq-remainder-bound} controls the first
through curvature and shorter-horizon conditional mean squared errors;
Equation \eqref{eq-propagated-bound} additionally accounts for
propagation through \(G_h\), defined in Appendix \ref{sec-proof_ma}. For
fixed \(L\geq0\), this factor} is \CPTSadd{bounded by \((1-L)^{-1}\) when \(L<1\),
equals \(h-1\) when \(L=1\), and grows geometrically with \(h\) when
\(L>1\).}

\CPTSadd{For the small-perturbation analysis, fix \(t,h,d\) with \(h\geq2\) and
consider} a \CPTSadd{family of processes indexed by \(\delta>0\)}, \CPTSadd{each satisfying}
the \CPTSadd{assumptions} for \CPTSadd{Equation \eqref{eq-ar}} in \CPTSadd{Proposition~\ref{prp-ar}.
Suppose that} the \CPTSadd{following bounds hold almost surely for finite
deterministic constants \(M,K,\sigma\geq0\) independent} of \CPTSadd{\(\delta\):}

\begin{enumerate}
\def\labelenumi{\arabic{enumi}.}
\tightlist
\item
  \CPTSadd{The Hessian bounds in Proposition~\ref{prp-ar} hold with
  \(M_{t,j}\leq M\)} for \CPTSadd{\(2\leq j\leq h\).}
\item
  \CPTSadd{The gradients along} the \CPTSadd{expansion segments satisfy} \[
  \CPTSadd{\sup_{u\in[0,1]}\left\|\nabla F_{t,j}\bigl(\bm m_{t,j}+u(\bm z_{t,j}-\bm m_{t,j})\bigr)\right\|_2\leq K,
  \qquad 2\leq j\leq h.}
  \]
\item
  \CPTSadd{The innovation conditional variances satisfy
  \(\mathbb E(\varepsilon_{t+j}^2|\mathcal F_{t+j-1})\leq\sigma^2\delta^2\)
  for \(1\leq j\leq h\).}
\end{enumerate}

\noindent \CPTSadd{Dependence on \(\delta\)} is \CPTSadd{suppressed in} the \CPTSadd{notation.}

\CPTSadd{By Cauchy--Schwarz}, \CPTSadd{\(\sum_{i=1}^{p_j}|\phi_{t,j,i}|\leq\sqrt d K\), so
\(L=\sqrt d K\) is admissible}.

\CPTSadd{Write \(V_j=\mathbb E(e_{t+j|t}^2|\mathcal F_t)\) for \(1\leq j\leq h\).
Conditional mean optimality gives, for \(j\geq2\),} \[
\CPTSadd{V_j\leq\mathbb E\!\left((y_{t+j}-F_{t,j}(\bm m_{t,j}))^2\,|\,\mathcal F_t\right).}
\] \CPTSadd{Since \(F_{t,j}(\bm z_{t,j})-F_{t,j}(\bm m_{t,j})\) is
\(\mathcal F_{t+j-1}\)-measurable, its cross term with
\(\varepsilon_{t+j}\) has conditional expectation zero. The gradient
bound gives
\(|F_{t,j}(\bm z_{t,j})-F_{t,j}(\bm m_{t,j})|\leq K\|\bm\Delta_{t,j}\|_2\),
using} the \CPTSadd{increment defined in Appendix \ref{sec-proof_ar}. Together
with the innovation variance bound}, \CPTSadd{this yields} \[
\CPTSadd{V_1\leq\sigma^2\delta^2,\qquad
V_j\leq\sigma^2\delta^2+K^2\sum_{i=1}^{p_j}V_{j-i},\qquad 2\leq j\leq h.}
\] \CPTSadd{Induction gives \(V_j\leq A_j\delta^2\),} where \CPTSadd{\(A_1=\sigma^2\) and
\(A_j=\sigma^2+K^2\sum_{i=1}^{p_j}A_{j-i}\). Taking
\(C=\max_{2\leq j\leq h}\sum_{i=1}^{p_j}A_{j-i}\), independent} of
\CPTSadd{\(\delta\), and substituting \(V=C\delta^2\) into} the \CPTSadd{remainder bounds
yields, almost surely,} \[
\CPTSadd{\mathbb E(|\widetilde R_{t,j}|\,|\,\mathcal F_t)\leq MC\delta^2,\qquad 2\leq j\leq h,
\qquad
\mathbb E(|U_{t,h}|\,|\,\mathcal F_t)\leq MCG_h\delta^2.}
\] \CPTSadd{Taking expectations shows} that\CPTSadd{, under} the \CPTSadd{stated conditions, both
remainders have mean absolute size \(O(\delta^2)\) and are negligible
relative to \(\delta\). Specifically,
\(\mathbb E|\widetilde R_{t,j}/\delta|\to0\) for \(2\leq j\leq h\) and
\(\mathbb E|U_{t,h}/\delta|\to0\) as \(\delta\downarrow0\).}

\CPTSadd{These results quantify linear approximation errors as the innovation
scale decreases at fixed} forecast \CPTSadd{horizons, rather than as the
training-sample size increases}. \CPTSadd{At a given perturbation scale, large
curvature or propagation factors may make the bounds uninformative.}

\subsection{\texorpdfstring{Proof of
Proposition~\ref{prp-cov_rt}}{Proof of Proposition~}}\label{sec-proof_cov_rt}

\begin{proof}
Let \CPTSadd{\(E_T=\sum_{i=h+1}^{T}(\mathrm{err}_{i|i-h}-\alpha)\)}. The
inequality \CPTSadd{in} Equation \eqref{eq-cov_rt} \CPTSadd{is equivalent to
\(|E_T|\leq cg(T-h)+h\)}. We \CPTSadd{first} prove the \CPTSadd{upper bound by strong
induction}, with the \CPTSadd{lower bound} following analogously.

For \CPTSadd{\(T=h+1,\dots,2h\)}, \CPTSadd{\(|E_T|\leq T-h\leq h\leq cg(T-h)+h\)}, \CPTSadd{since}
\(g\) is nonnegative and \CPTSadd{\(0\leq\mathrm{err}_{i|i-h}\leq1\)}. Thus \CPTSadd{both
bounds hold over the initial range}.

Now \CPTSadd{let \(T\geq2h+1\) and assume the upper bound holds at all earlier
indices}. \CPTSadd{The threshold evaluated at time \(T\) is} \[
\CPTSadd{q_{T|T-h}=r_{T-h}(E_{T-h}).}
\] We partition the argument into \CPTSadd{two} cases\CPTSadd{. If
\(E_{T-h}\geq cg(T-2h)\)}, Equation \eqref{eq-saturation_h} \CPTSadd{gives
\(q_{T|T-h}\geq b\)}, \CPTSadd{so \(\mathrm{err}_{T|T-h}=0\). By the induction
hypothesis} and \CPTSadd{monotonicity of \(g\),} \[
\CPTSadd{E_T=E_{T-1}-\alpha\leq cg(T-1-h)+h\leq cg(T-h)+h.}
\] \CPTSadd{Otherwise, using \(\mathrm{err}_{i|i-h}-\alpha\leq1-\alpha\) over the
last \(h\) observations gives} \[
\CPTSadd{E_T=E_{T-h}+\sum_{i=T-h+1}^{T}(\mathrm{err}_{i|i-h}-\alpha)
<cg(T-2h)+h(1-\alpha)\leq cg(T-h)+h.}
\] \CPTSadd{This completes the induction for the upper bound}. \CPTSadd{For the lower
bound}, \CPTSadd{the same induction applies: if \(E_{T-h}\leq-cg(T-2h)\),
saturation and \(s_{T|T-h}>-b\) imply \(\mathrm{err}_{T|T-h}=1\), so} \[
\CPTSadd{E_T=E_{T-1}+1-\alpha\geq-cg(T-1-h)-h\geq-cg(T-h)-h.}
\] \CPTSadd{Otherwise, \(\mathrm{err}_{i|i-h}-\alpha\geq-\alpha\) yields} \[
\CPTSadd{E_T>-cg(T-2h)-h\alpha\geq-cg(T-h)-h.}
\] \CPTSadd{The argument also applies when \(b=\infty\) because every score} is
\CPTSadd{finite}. \CPTSadd{Dividing by \(T-h\) proves Equation \eqref{eq-cov_rt}.}

\CPTSadd{Regarding the long-run conclusion}, \CPTSadd{\(g(T-h)/(T-h)\to0\) by sublinearity,}
and \CPTSadd{\(h/(T-h)\to0\) for fixed \(h\)}. \CPTSadd{Hence the empirical tail
miscoverage rate converges to \(\alpha\).}
\end{proof}

\subsection{\texorpdfstring{Proof of
Corollary~\ref{cor-cov_qt}}{Proof of Corollary~}}\label{sec-proof_cov_qt}

\begin{proof}
\CPTSadd{Unraveling the recursion gives} \[
\CPTSadd{q_{t+h|t}=q_{2h|h}+\eta\sum_{i=h+1}^{t}(\mathrm{err}_{i|i-h}-\alpha).}
\] \CPTSadd{Thus}, \CPTSadd{take \(r_t(x)=q_{2h|h}+\eta x\)}, \CPTSadd{\(g\equiv1\)}, \CPTSadd{and
\(c=(b+|q_{2h|h}|)/\eta\)}. \CPTSadd{If \(x\geq c\)}, \CPTSadd{then \(r_t(x)\geq b\), and if
\(x\leq-c\), then \(r_t(x)\leq-b\), so Equation \eqref{eq-saturation_h}
holds. Applying Proposition~\ref{prp-cov_rt} gives} the \CPTSadd{stated bound and
long-run calibration}.
\end{proof}

\subsection{\texorpdfstring{Proof of
Corollary~\ref{cor-cov_acmcp}}{Proof of Corollary~}}\label{sec-proof_cov_acmcp}

\begin{proof}
\CPTSadd{Let \(s_{t+h|t}^{*}=s_{t+h|t}-\hat q_{t+h|t}\) and
\(q_{t+h|t}^{*}=q_{t+h|t}-\hat q_{t+h|t}\). Since the scores and
baseline values are finite}, \CPTSadd{these transformations are well-defined even
when the thresholds are infinite}, and \CPTSadd{\(s_{t+h|t}^{*}\in(-b,b)\) by
assumption}. \CPTSadd{Subtracting the same finite baseline preserves the tail
miscoverage indicator:
\(\mathrm{err}_{t+h|t}=\mathbbm{1}\{s_{t+h|t}^{*}>q_{t+h|t}^{*}\}\)}.
\CPTSadd{Equation \eqref{eq-acmcp_1} therefore becomes} \[
\CPTSadd{q_{t+h|t}^{*}=r_t\left(\sum_{i=h+1}^{t}\bigl(\mathbbm{1}\{s_{i|i-h}^{*}>q_{i|i-h}^{*}\}-\alpha\bigr)\right).}
\] \CPTSadd{The transformed recursion satisfies the assumptions of
Proposition~\ref{prp-cov_rt}} with \CPTSadd{the same \(b\)}, \CPTSadd{\(c\), and \(g\).
Applying} that \CPTSadd{proposition} gives the \CPTSadd{stated finite-sample bound and
long-run calibration for the unchanged indicators}.
\end{proof}

\subsection{\texorpdfstring{Proof of
Corollary~\ref{cor-cov-two-sided}}{Proof of Corollary~}}\label{sec-proof_cov_two_sided}

\begin{proof}
\CPTSadd{Common saturation constants} can \CPTSadd{be chosen by taking} the \CPTSadd{maximum of the
two constants and the pointwise maximum of the two admissible functions.
Apply Corollary~\ref{cor-cov_acmcp} separately to the lower- and
upper-tail recursions, replacing its generic target \(\alpha\) by
\(\alpha^{-}\) and \(\alpha^{+}\), respectively. With \(n=T-h\), this
gives} \[
\CPTSadd{\left|\frac{1}{n}\sum_{t=h+1}^{T}\mathrm{err}_{t|t-h}^{a}-\alpha^{a}\right|\leq \frac{c\cdot g(n)+h}{n},
\qquad a\in\{-,+\}.}
\] \CPTSadd{The interval fails to cover \(y_t\) if and only if} at \CPTSadd{least one tail
constraint fails, including when the interval is empty}. \CPTSadd{Thus}, \[
\CPTSadd{\mathbbm{1}\{y_t\notin\hat{\mathcal C}_{t|t-h}^{\alpha}\}
=\mathrm{err}_{t|t-h}^{-}+\mathrm{err}_{t|t-h}^{+}-\mathrm{err}_{t|t-h}^{-}\mathrm{err}_{t|t-h}^{+}.}
\] \CPTSadd{Dropping the nonnegative overlap term}, \CPTSadd{summing,} and using
\CPTSadd{\(\alpha^{-}+\alpha^{+}=\alpha\) yields} \[
\CPTSadd{\frac{1}{n}\sum_{t=h+1}^{T}\mathbbm{1}\{y_t\notin\hat{\mathcal C}_{t|t-h}^{\alpha}\}
\leq \alpha+2\frac{c\cdot g(n)+h}{n},}
\] \CPTSadd{which is equivalent to Equation \eqref{eq-two-sided-finite}. Because
\(g\) is sublinear and \(h\) is fixed},
\CPTSadd{\(\bigl(c\cdot g(n)+h\bigr)/n\to0\), giving the long-run coverage lower
bound.}

\CPTSadd{For the absolute deviation bound, simultaneous violations of both tail
bounds imply} \[
\CPTSadd{q_{t|t-h}^{+}<y_t-\hat y_{t|t-h}< -q_{t|t-h}^{-},}
\] \CPTSadd{so \(\hat y_{t|t-h}-q_{t|t-h}^{-}>\hat y_{t|t-h}+q_{t|t-h}^{+}\)}, \CPTSadd{a
crossing of} the \CPTSadd{displayed endpoints}. This \CPTSadd{comparison remains
well-defined} for \CPTSadd{infinite thresholds. Thus} the \CPTSadd{average} of \CPTSadd{the products
of the two tail indicators is at most \(\chi_{T,h}\). The union identity
and the triangle inequality give} \[
\CPTSadd{\begin{aligned}
\left|\frac1n\sum_{t=h+1}^{T}\mathbbm{1}\{y_t\in\hat{\mathcal C}_{t|t-h}^{\alpha}\}-(1-\alpha)\right|
&\leq\sum_{a\in\{-,+\}}\left|\frac1n\sum_{t=h+1}^{T}\mathrm{err}_{t|t-h}^{a}-\alpha^a\right|
+\frac1n\sum_{t=h+1}^{T}\mathrm{err}_{t|t-h}^{-}\mathrm{err}_{t|t-h}^{+}\\
&\leq2\frac{c\cdot g(n)+h}{n}+\chi_{T,h}.
\end{aligned}}
\] \CPTSadd{This proves Equation \eqref{eq-two-sided-enhanced} and exact long-run
calibration when \(\chi_{T,h}\to0\).}

\CPTSadd{For the iteration in} Proposition~\ref{prp-cov_rt}\CPTSadd{, the same proof
applies with zero baselines}.
\end{proof}

\subsection{\CPTSadd{Two-sided coverage for constant-step quantile
tracking}}\label{sec-two-sided-constant-step}

\begin{corollary}[]\protect\hypertarget{cor-cov-two-sided-qt}{}\label{cor-cov-two-sided-qt}

\CPTSadd{Fix \(h\in[H]\) and targets \(\alpha^{-},\alpha^{+}>0\) with
\(\alpha^{-}+\alpha^{+}=\alpha\in(0,1)\). Apply} the \CPTSadd{iteration in
Corollary~\ref{cor-cov_qt} separately} to the \CPTSadd{two signed tail scores and
construct the interval in Equation \eqref{eq-two-sided-pi}. Suppose each
tail \(a\in\{-,+\}\) satisfies the assumptions} of
\CPTSadd{Corollary~\ref{cor-cov_qt} with finite score bound \(b^a\), step size
\(\eta^a>0\), and finite initial threshold \(q^a_{2h|h}\). Then, for
every \(T\geq h+1\),}
\begin{equation}\CPTSadd{\phantomsection\label{eq-two-sided-constant-step}{
\frac{1}{T-h}\sum_{t=h+1}^{T}\mathbbm{1}\{y_t\in\hat{\mathcal C}_{t|t-h}^{\alpha}\}
\geq 1-\alpha-\sum_{a\in\{-,+\}}\frac{b^a+|q^a_{2h|h}|+\eta^a h}{\eta^a(T-h)}.
}}\end{equation} \CPTSadd{The absolute deviation of} the \CPTSadd{empirical coverage from
\(1-\alpha\) is bounded by the sum on the right plus the
endpoint-crossing frequency \(\chi_{T,h}\) defined in
Corollary~\ref{cor-cov-two-sided}}. \CPTSadd{Consequently}, the \CPTSadd{limit inferior of
the empirical coverage is at least \(1-\alpha\), and the coverage
converges to \(1-\alpha\) if \(\chi_{T,h}\to0\)}.

\end{corollary}

\begin{proof}
\CPTSadd{With \(n=T-h\)}, \CPTSadd{Corollary~\ref{cor-cov_qt} gives} \[
\CPTSadd{\left|\frac1n\sum_{t=h+1}^{T}\mathrm{err}_{t|t-h}^{a}-\alpha^a\right|
\leq\frac{b^a+|q^a_{2h|h}|+\eta^a h}{\eta^a n},\qquad a\in\{-,+\}.}
\] \CPTSadd{Applying} the \CPTSadd{union identity and endpoint-crossing argument in
Appendix \ref{sec-proof_cov_two_sided} gives Equation
\eqref{eq-two-sided-constant-step} and the absolute deviation bound. For
fixed \(h\), the two tail bounds tend} to \CPTSadd{zero, giving the stated
long-run conclusions}.
\end{proof}

\subsection{\CPTSadd{Local coverage bounds}}\label{sec-local-coverage}

\CPTSadd{Fix \(h\in[H]\)} and \CPTSadd{a window \(\{a+1,\ldots,a+m\}\), where \(a\geq h\)
and \(m\geq1\) are integers. Under} the \CPTSadd{assumptions of
Corollary~\ref{cor-cov_acmcp}}, \CPTSadd{set \(D_h=0\) and \(D_u=c\cdot g(u-h)+h\)}
for \CPTSadd{\(u\geq h+1\)}. \CPTSadd{Differencing the cumulative deviations at \(a+m\) and
\(a\) and applying the triangle inequality yields}
\begin{equation}\CPTSadd{\phantomsection\label{eq-window-tail}{
\left|\frac1m\sum_{t=a+1}^{a+m}(\mathrm{err}_{t|t-h}-\alpha)\right|
\leq\frac{D_{a+m}+D_a}{m}.
}}\end{equation} \CPTSadd{Under the assumptions of
Corollary}~\CPTSadd{\ref{cor-cov-two-sided}, with common \(c\)} and \CPTSadd{\(g\) for both
tails and \(\alpha=\alpha^{-}+\alpha^{+}\),} the \CPTSadd{union bound gives}
\begin{equation}\CPTSadd{\phantomsection\label{eq-window-coverage}{
\frac1m\sum_{t=a+1}^{a+m}\mathbbm{1}\{y_t\in\hat{\mathcal C}_{t|t-h}^{\alpha}\}
\geq1-\alpha-2\frac{D_{a+m}+D_a}{m}.
}}\end{equation} \CPTSadd{If \((D_{a+m}+D_a)/m\to0\) along a sequence of windows,
empirical window coverage has limit inferior at least \(1-\alpha\) and
converges to \(1-\alpha\) if the within-window endpoint-crossing
frequency also vanishes, by Appendix \ref{sec-proof_cov_two_sided}}.
\CPTSadd{These bounds may be conservative and uninformative for fixed-length
windows.}

\subsection{\CPTSadd{Sensitivity to uncertainty in future
predictors}}\label{sec-predictor-uncertainty}

\CPTSadd{We assess the effect of unknown future predictors using} \[
\CPTSadd{x_t=0.7x_{t-1}+u_t,\qquad
y_t=0.8y_{t-1}-0.5y_{t-2}+x_t+\varepsilon_t,}
\] \CPTSadd{where \(u_t\overset{\mathrm{iid}}{\sim}N(0,\tau^2)\) and
\(\varepsilon_t\overset{\mathrm{iid}}{\sim}N(0,1)\) are mutually
independent}, \CPTSadd{with \(\tau\in\{0.25,1\}\). Each of 100 replications
retains 5000 observations after a 500-observation burn-in. At each
forecast origin}, an \CPTSadd{ARX(2) model} with \CPTSadd{an intercept is estimated by least
squares over a rolling window of 500 observations. We compare forecasts
using observed future predictors (Known) with forecasts using
\(\hat x_{t+j|t}=0.7^j x_t\) (Forecast). The autoregressive coefficient
of \(x_t\) is treated as known} to \CPTSadd{isolate future predictor uncertainty
from its estimation error}.

\CPTSadd{Both settings use the same simulated paths and the AcMCP settings in
Section~\ref{sec-simar}}, \CPTSadd{but are calibrated separately. We evaluate
horizons \(h=1,2,3\) at a target coverage of \(90\%\). Local coverage
MAE measures} the \CPTSadd{average absolute deviation of 500-observation rolling
coverage from \(90\%\).}

\CPTSadd{Table~\ref{tbl-predictor} shows that both settings maintain coverage
close to \(90\%\)}, with \CPTSadd{local coverage MAE below \(0.3\) percentage
points. Interval widths} in \CPTSadd{the Known setting are nearly unchanged across
the two values of \(\tau\)}, \CPTSadd{whereas forecasting future predictors leads
to wider intervals, with larger increases at higher \(\tau\) and longer
horizons. At \(h=3\),} the \CPTSadd{increase is about \(8\%\) for \(\tau=0.25\)}
and \CPTSadd{\(97\%\) for \(\tau=1\)}. Thus, \CPTSadd{in this simulation,} AcMCP
\CPTSadd{accommodates predictor uncertainty primarily through wider intervals
while maintaining empirical coverage}.

\begin{table}

\caption{\label{tbl-predictor}\CPTSadd{AcMCP with observed (Known) and forecast (Forecast) future predictors. Entries are means over 100 paired simulation replications, with Monte Carlo standard errors in parentheses. pp denotes percentage points.}}

\centering{

\resizebox{\linewidth}{!}{
\begin{tabular}{lrrrrrr}
\toprule
 & \multicolumn{2}{c}{Coverage (\%)} & \multicolumn{2}{c}{Mean interval width} & \multicolumn{2}{c}{Local coverage MAE (pp)} \\
\cmidrule(l{3pt}r{3pt}){2-3} \cmidrule(l{3pt}r{3pt}){4-5} \cmidrule(l{3pt}r{3pt}){6-7}
 & \multicolumn{1}{c}{Known} & \multicolumn{1}{c}{Forecast} & \multicolumn{1}{c}{Known} & \multicolumn{1}{c}{Forecast} & \multicolumn{1}{c}{Known} & \multicolumn{1}{c}{Forecast} \\
\midrule
 & \multicolumn{6}{c}{$\tau=0.25$} \\[0.3em]
$h=1$ & 89.99 (0.003) & 89.99 (0.004) & 3.55 (0.005) & 3.65 (0.005) & 0.20 (0.001) & 0.19 (0.001) \\
$h=2$ & 89.99 (0.004) & 89.99 (0.004) & 4.72 (0.008) & 5.00 (0.010) & 0.24 (0.002) & 0.24 (0.002) \\
$h=3$ & 89.99 (0.004) & 89.99 (0.005) & 4.85 (0.009) & 5.26 (0.011) & 0.27 (0.002) & 0.27 (0.002) \\
\cmidrule(l{3pt}r{3pt}){2-7}
 & \multicolumn{6}{c}{$\tau=1$} \\[0.3em]
$h=1$ & 89.99 (0.003) & 89.99 (0.003) & 3.56 (0.005) & 4.99 (0.007) & 0.20 (0.001) & 0.19 (0.001) \\
$h=2$ & 89.99 (0.004) & 89.99 (0.004) & 4.72 (0.009) & 8.10 (0.016) & 0.24 (0.002) & 0.24 (0.001) \\
$h=3$ & 89.98 (0.005) & 89.99 (0.005) & 4.85 (0.009) & 9.57 (0.021) & 0.27 (0.002) & 0.29 (0.002) \\
\bottomrule
\end{tabular}}

}

\end{table}%





\end{document}
